%% AER submission format: 12pt, double-spaced, 1-inch margins,
%% page numbers, article class (AER-compatible).
\documentclass[12pt]{article}

\usepackage{amsmath,amssymb,amsfonts}
\usepackage[T1]{fontenc}
\usepackage[utf8]{inputenc}
\usepackage{graphicx}
\usepackage{booktabs}
\usepackage{threeparttable}
\usepackage[letterpaper,margin=1in]{geometry}
\usepackage{setspace}
\usepackage{appendix}
\usepackage[hyphens]{url}
\usepackage{xurl}
\usepackage{hyperref}
\hypersetup{colorlinks=true, linkcolor=blue, citecolor=blue,
            filecolor=blue, urlcolor=blue}
\usepackage{microtype}
\usepackage{natbib}
\usepackage{adjustbox}
\usepackage{multirow}
\usepackage{lmodern}
\usepackage{caption}
\usepackage{subfig}

\setlength{\emergencystretch}{3em}
\setlength{\bibsep}{0pt plus 0.3ex}
\doublespacing
\graphicspath{{figures/}{../output/figures/}}
\pdfsuppresswarningpagegroup=1

% AER-style captions: small, single-spaced, hang indent
\captionsetup{font=small, labelfont=bf, labelsep=period}
\captionsetup[table]{singlelinecheck=off}

% AER title formatting
\title{\textbf{Set-Asides Reshape Bids, Not Bidders:\\
Evidence from Brazilian Medical Procurement}%
\thanks{\emph{This version: April 2026 (v2-structural).}
A parallel paper develops the reduced-form evidence in
Genicolo-Martins (2026a). Replication materials are documented
in Appendix~L and available at
\url{https://github.com/darciogm/paper2-me-epp/v2-structural}.
All errors are my own.}}

\author{Darcio Genicolo-Martins\thanks{INSPER, São Paulo, Brazil.
Email: \href{mailto:darciogm1@insper.edu.br}{darciogm1@insper.edu.br}.}}

\date{April 2026}

\newcommand{\citetp}[1]{\citet{#1}}

\begin{document}

\maketitle
\thispagestyle{empty}

\begin{center}
\begin{minipage}{0.85\textwidth}
\small
\begin{singlespace}
\textbf{Abstract.}
Does a small-firm set-aside raise prices by excluding competitors,
or by changing how the remaining firms bid? For Brazil's March~2018
extension of the SME-only rule to medical supplies, the
\citet{campo2003} estimator on 33{,}706 auctions decomposes the
10--11\% price effect: \textbf{66\% is within-auction bidding
response, 34\% is cost-primitive entry.} A testable
primitive-invariance condition for non-SMEs holds tightly
(KS $D = 0.03$). Participation shifts are large ($-62\%$/$+52\%$),
but the price effect lives inside the auction, not at the door.
Deadweight loss is 11.8\% of procurement value; a top-quartile
exemption recovers 95\% of the fiscal cost. \textbf{Set-asides
reshape bids, not bidders.}

\vspace{0.5em}

\noindent \textbf{JEL codes:} H32, H57, L26, L53, D44, C14.

\noindent \textbf{Keywords:} public procurement; small-firm
set-asides; structural auction estimation; asymmetric GPV; welfare
counterfactuals.
\end{singlespace}
\end{minipage}
\end{center}

\clearpage
\setcounter{page}{1}
\pagestyle{plain}

\clearpage
\section{Introduction}

When a government restricts public procurement to small firms, does the
price premium come from the firms it excludes, or from the firms it
keeps? The question is old---set-aside rules have been on the books in
developed and developing economies since at least the 1950s---and the
answer determines what reform looks like. If the premium runs through
the extensive margin (fewer firms bid), the policy fix is to expand the
small-firm roster: subsidies, registration relief, capacity-building.
If instead it runs through the intensive margin (the same firms bid less
aggressively because they no longer face large competitors), no roster
expansion will close the gap---the only lever is to shrink the
\emph{scope} of the rule itself. Policy prescriptions that differ by an
order of magnitude turn on a quantity the literature has mostly been
unable to measure directly.

This paper measures it, using a quasi-experiment in Brazilian medical
procurement. Between August 2014 and February 2018, medical supplies
(Group 65 in the BEC product catalog) were the \emph{only} major
product group in S\~ao Paulo state that remained open to general
competition; 76 other product groups had been under the SME-only
tendering rule since the federal 2014 mandate took effect. The
exception was the outcome of a joint agreement between the state
government and its audit court (TCE-SP) that medical supplies required
open competition on public-health grounds. That exception ended in
March 2018, when a legal opinion from a separate state agency
(the \emph{Procuradoria-Geral do Estado}, PGE-SP) held the exemption
unconstitutional on isonomy grounds: medical supplies could not be
procedurally distinguished from other groups. The SME-only rule
extended to Group 65 overnight. The reversal was triggered by a
bureaucratic legal argument, not by any procurement-market outcome; it
is plausibly exogenous to anything happening in Group-65 bidding.

\paragraph{The approach.}
Section~\ref{sec:rf} establishes via difference-in-differences
that the policy raises Group-65 procurement prices by 10--11\%
over the 18-month window, a figure robust to the
\citet{callaway2021}, \citet{sun2021}, and \citet{goodmanbacon2021}
estimators. Reduced-form mediation, however, cannot separate the
two channels of interest. A \citet{gelbach2016} decomposition locates roughly 95\%
of the effect outside the observable entry and winner-composition
channels. That residual is what the present paper explains: using bid
distributions recovered from the BEC electronic platform, I estimate
an asymmetric first-price sealed-bid (FPSB) model on the \emph{Convite}
modality of Group-65 procurement (33,706 auctions with at least two
bidders; 24,346 bidder-auctions entering the CPV inversion) via the
\citet{campo2003} extension of \citet{gpv2000}. The estimation
delivers type-specific cost distributions
$F_c^{\text{SME}}(\cdot)$ and $F_c^{\text{NonSME}}(\cdot)$.
Because the regime switch affects \emph{which} firms bid but not
their primitive costs, the non-SME cost distribution should be
invariant across the break. That invariance, when it holds, both
validates the model and pins down the channel structure.

\paragraph{Three findings.}

\emph{First}, the identifying restriction holds where the model
predicts it should. In Convite, where non-SMEs continue bidding in
many auctions both before and after the cutoff, their recovered cost
distribution is statistically invariant across the March-2018 break:
mean shift of 0.001 of reference price, KS $D = 0.03$. Under the
Preg\~ao Haile-Tamer upper bound, the pattern flips---it is the SME
bound that is nearly invariant ($D = 0.03$), while the non-SME bound
drifts upward by 0.04 of reference. The drift reflects selection among
the shrinking pool of non-SMEs who keep bidding despite near-zero
post-policy win probability, not a movement in primitive costs. In
each modality, the primitive that survives the test is the one whose
bidder pool was not re-selected by the policy.

\emph{Second}, the SME cost distribution does shift upward---by 7.7\%
relative to its Pre-period mean, equivalent to 4 percentage points
in ref-normalized units (KS $D = 0.08$, $p < 10^{-10}$). This is the
compositional footprint of the policy: when Group 65 becomes
SME-only, marginal higher-cost small firms---firms that previously
specialized in other product groups or new entrants---enter the
Group-65 candidate pool. The SME primitive is no longer the same pool
before and after. This ``entry margin'' is, to my knowledge, the first
structural measurement of pool re-composition under an SME set-aside.

\emph{Third}, the structural price effect decomposes as
\textbf{$\sim$66\% intensive margin and $\sim$34\% entry margin}.
The intensive share is 74\% at $N{=}2$, 47\% at $N{=}3$, 70\% at
$N{=}4$, and 72\% at $N{\geq}5$---a rough two-thirds across auction
sizes. This is the central structural finding. It does the work
the Gelbach residual could not: the 95\% unexplained by observable
entry and winner composition is, at its core, \emph{within-auction
bidding behavior}---the same SMEs bidding differently once their
larger competitors are gone---not another observable margin waiting
to be added. Participation responds too, and substantially: the
reduced-form entry regression shows the non-SME entry rate falls
62\% under the restriction ($t = 43$) while the SME entry rate
rises 52\% ($t = 38$). These shifts amplify the total welfare cost
the policy generates, but they are not where the price effect
primarily lives. The roughly R\$30--33 (US\$9--10) type-symmetric
entry cost implied by the free-entry zero-profit condition means
that a
rigorous \citet{atheyseira2011} nested entry-bidding model---
future work---should refine magnitudes without overturning the
intensive-margin primacy of the fixed-pool result.

\paragraph{Welfare counterfactual.}
Using the recovered cost distributions, I simulate equilibrium
procurement prices under a \emph{value-threshold exemption}---
applying the SME-only rule only below a reference-price cutoff.
The structural Pareto frontier is strikingly concave: exempting
the top quartile of items by reference price preserves the SME
preference on 75\% of items \emph{by count} while recovering
\textbf{95\%} of the fiscal cost of the uniform rule. This reform
is not just arithmetic---the structural model shows that the
top-quartile items are where the intensive margin bites hardest,
so reopening them is also where cost recovery is largest.

\paragraph{Contributions.}
The findings above distill into three methodological contributions
to the literature:

\noindent\textbf{(i) First structural decomposition of an SME
set-aside into intensive and entry channels.} The public-procurement
literature has either measured SME effects in reduced
form~\citep{marion2007, denes1997, ferraz2015} or modeled bid
preferences structurally without the clean
\emph{restriction}-versus-\emph{preference}
contrast~\citep{krasnokutskayaseim2011, athey2013}. I combine the
two: reduced-form identification via a discrete legal cutoff,
followed by structural estimation of type-specific cost primitives.
The headline finding---that within-auction bidding carries the
majority of the policy's price effect---is directly relevant to
the design of SME preference regimes in Brazil, the EU, and the
130 other jurisdictions documented in \citet{oecd2018, sigma2016,
hoekman2020}.

\noindent\textbf{(ii) A selection-conditional consistency check on
primitives.} In the Convite CPV estimator, the non-SME cost
distribution is tightly identified across the regime break: it
enters a (same) set of auctions with (essentially) the same
bidder-pool structure before and after the policy. The resulting
KS $D = 0.03$ (mean shift $0.001$) is a consistency check on the
model under the identifying assumption that the policy did not
re-select which non-SMEs bid in the Convite Pre vs Post samples.
This condition plausibly holds in Convite because non-SMEs remain
eligible to bid in Convite auctions not subject to the SME-only
restriction, and the regime change affects a subset of Group-65
auctions rather than the entire non-SME choice set. Under Preg\~ao,
where the Haile-Tamer upper bound is applied to final bids and
non-SMEs continue to bid despite near-zero post-policy win
probability, the upper bound on the non-SME cost distribution does
shift ($D = 0.06$), consistent with selection of the remaining
non-SME bidders toward higher-cost types. The paper's
identification relies on the Convite cross-regime check, with the
Preg\~ao result as a directional robustness.

\noindent\textbf{(iii) Welfare-based structural evaluation of a
value-threshold reform.} I simulate a concrete policy alternative
---exempting high-reference-price items from the SME-only rule
while keeping the preference on the rest---and show that the
structural primitives deliver a Pareto frontier considerably more
concave than an accounting-identity calculation would suggest.
The top-quartile exemption recovers 95\% of the fiscal cost on
25\% of items, grounded in asymmetric-bidder welfare accounting
rather than mechanical extrapolation.

\paragraph{Relation to the literature.}
The paper connects two strands. The empirical SME procurement
literature---\citet{marion2007, denes1997, ferraz2015, szucs2024}---documents
price and efficiency effects of set-aside and preference rules in
reduced form. The structural auction literature---\citet{gpv2000,
campo2003, krasnokutskaya2011, athey2013, bajari2014,
bajarihortacsu2005, hendricks2003}---recovers cost primitives from
bid data and simulates counterfactuals, typically in settings without
a clean regime break. This paper bridges the two:
the March 2018 break provides the causal cut; CPV asymmetric estimation
recovers the primitives; and the structural counterfactuals quantify
welfare. The approach is closest in spirit to
\citet{krasnokutskayaseim2011} on California highway procurement with
bid preferences, but with three distinctions: (i) the regime contrast
is \emph{restriction} vs \emph{open}, not a 5\% preference; (ii) the
asymmetric bidder-type framework is \citet{campo2003} rather than
symmetric GPV with unobserved heterogeneity as in
\citet{krasnokutskaya2011}; and (iii) the March 2018 regime break
provides a \emph{testable} primitive-invariance condition on the
non-SME cost distribution, which \citet{krasnokutskayaseim2011} assume
rather than test.

For comparison of magnitudes: \citet{krasnokutskayaseim2011} analyze
697 Caltrans projects with 3{,}034 bids, an average of 4.4 bidders per
project (1.75 small + 2.62 large). Their counterfactual estimates the
California 5\% Small Business Preference raises procurement cost by
0.5\% and shifts small firms' share of work from 12.5\% to 15.6\%; a
25\% bid discount would raise costs by 4.1\% (their Table 11). The
present paper's set-aside regime (complete exclusion of non-SMEs) is
an order-of-magnitude stronger policy and generates a proportionally
larger effect: a 10.4\% fiscal cost, 11.8\% DWL of procurement value,
roughly 20$\times$ the \citet{krasnokutskayaseim2011} 5\% baseline.
That ratio is broadly consistent with moving from a mild preference
to a full restriction. Type-specific cost asymmetries are broadly consistent:
\citet{krasnokutskayaseim2011} estimate small firms bid 5.4\%
(parametric) to 8.1\% (OLS) higher than large firms; we estimate
small firms bid about 17\% higher under the Preg\~ao Haile-Tamer
upper bound (H1-only; one-sided bound), about 4.7\% higher under the
canonical Convite CPV (Appendix~C), and essentially zero under the
proxy-classification Convite CPV (where the proxy pools firms by
registry class and the cross-type distinction is noisier at the
margin).

\paragraph{Roadmap.}
Section~\ref{sec:inst} describes the institutional background---the
Brazilian procurement regime, the BEC platform, and the 2014--2018
regulatory history of Group 65 that motivates the identifying cut.
Section~\ref{sec:rf} summarizes the reduced-form baseline
establishing the 10--11\% headline price effect and the 95\%
Gelbach residual that the structural exercise will decompose.
Section~\ref{sec:model} presents the asymmetric FPSB procurement model.
Section~\ref{sec:identification} states the identification strategy and
the Campo-Perrigne-Vuong inversion. Section~\ref{sec:data} describes
the bid-level data and the Convite modality.
Section~\ref{sec:estimation} estimates the model and reports the
primitive-invariance tests. Section~\ref{sec:decomposition}
delivers the headline 66/34 decomposition of the structural price
effect into intensive and entry channels.
Section~\ref{sec:endogenous} documents the participation response
(62\% non-SME exit, 52\% SME entry) and shows why it does not
overturn the within-auction primacy of the intensive margin.
Section~\ref{sec:welfare} presents the welfare analysis---accounting
identity, DWL magnitudes, annualization, and the value-threshold
counterfactual.
Section~\ref{sec:robustness} discusses robustness (Preg\~ao type-split
Haile-Tamer bounds, canonical SME, IPV-vs-CV diagnostic).
Section~\ref{sec:conclusion} concludes.

%% ============================================================================
\section{Institutional Background}
\label{sec:inst}

The identifying variation in this paper starts with a single
bureaucratic document. In early 2018, São Paulo's state attorney's
office (PGE-SP) issued an internal legal opinion that, on grounds
entirely unrelated to procurement outcomes, extended Brazil's
SME-only default to its last open product category: medical and
hospital supplies. A parecer on equal treatment took effect the
following month, and the auctions in this paper look different on
either side of that line. To see why the line is clean, it helps
to understand how the surrounding system works---the statutory
layer-cake of Brazilian procurement, the BEC platform that
operationalizes it in São Paulo, the 2014 federal mandate that
made SME-only tendering binding, the peculiar 2014--2018 exception
that kept Group 65 outside the rule, and the 2018 reinterpretation
that swept it in.

\paragraph{Brazilian public procurement: statutory framework.}
Brazilian procurement is a layered structure, and each layer leaves
a trace in the data. The foundational statute is Lei~8.666/1993,
which at the time of this study defined five classical modalities
sorted by value and formality: \emph{Concorr\^encia} for the
largest projects, \emph{Tomada de Pre\c{c}os} for pre-registered
suppliers, \emph{Convite} for invited sealed-bid auctions,
\emph{Concurso} for creative work, and \emph{Leil\~ao} for the
sale of public assets. Lei~10.520/2002 added a sixth modality,
the \emph{Preg\~ao}---a reverse auction that opens with sealed
proposals and closes with an iterative open-outcry phase---and
over time made it the workhorse of standardized procurement.
Convite and Preg\~ao are the two modalities this paper uses:
Convite because it is a first-price sealed-bid auction, exactly
the object the CPV estimator was designed for, and Preg\~ao
because its iterative structure puts it in the Haile-Tamer world
of ascending-bid bounds. Two additional pieces of legislation
matter for timing: Decreto~9.412/2018 updated the modality
ceilings mid-sample (raising the Convite non-engineering item
limit from R\$80{,}000 to R\$176{,}000), and Lei~14.133/2021 is
currently phasing in a unified replacement for 8.666 and 10.520.
Neither change affects the identifying window in this paper, but
both matter for external validity.

Layered on top of these modalities is a separate preferences regime
for small firms, introduced by Lei Complementar 123/2006---the
Microenterprise Statute, known in practice as LC~123. LC~123
defines two eligible firm sizes---\emph{Micro-empresa} (ME: annual
revenue $\leq$ R\$360{,}000) and \emph{Empresa de Pequeno Porte}
(EPP: $\leq$ R\$4.8M)---and grants them three cumulative
preferences: exclusive eligibility for items below R\$80{,}000
(Art.~48~I), subcontracting set-asides (Art.~48~II), and a 10\%
tie-breaking margin, reduced to 5\% in Preg\~ao (Art.~44~\S2).
For the first eight years of the statute these preferences were
optional tools that purchasing units could invoke or ignore.
That changed in 2014: LC~147/2014 converted them into binding
defaults, requiring any buyer who ran a general (non-restricted)
tender for a sub-R\$80{,}000 item to file a written justification
with the state audit court. The 2014 amendment is the proximate
source of the SME-only regime this paper studies---and the reason
the Group-65 exception that followed was such an anomaly.

\paragraph{BEC: the S\~ao Paulo procurement platform.}
The statutes describe rules; BEC enforces them. Since Decreto
estadual 47.297/2002, every São Paulo state purchase passes
through the \emph{Bolsa Eletr\^onica de Compras}, an electronic
procurement platform administered by the state treasury secretariat
(SEFAZ/SP). BEC is the single point of entry for 1{,}344 purchasing
units---state bureaus, participating municipalities, and accredited
private entities subject to state rules---and in 2019 it moved
roughly R\$13 billion (US\$3.7 billion) through its order book.
Cumulatively since 2005, BEC has handled more than R\$105 billion
(US\$30 billion), spread across 860{,}000 purchase offers and
4.8 million item-level transactions. That volume alone makes the
platform an unusually dense laboratory for procurement research.

Three nested registries structure the platform. Suppliers enroll
in the Cadastro Unificado de Fornecedores (CAUFESP), which
certifies their tax regularity, firm size, and procurement
eligibility; goods and services are cataloged in a three-tier
CADMAT taxonomy (\emph{group} $\to$ \emph{class} $\to$ \emph{item});
and each auction runs under one of the statutory modalities, with
Preg\~ao absorbing roughly 70\% of value, Convite 15\%, and
Dispensa eletr\^onica (below-R\$17{,}600 direct-purchase auctions)
15\%. In scope BEC parallels the federal ComprasNet, but it has
its own governance, its own supplier certification, and a different
closing protocol---BEC's Preg\~ao uses a soft-close rule rather
than the ``tempo aleat\'orio'' randomization that ComprasNet
adopted. For this paper the relevant CADMAT address is Group 65
(medical and hospital supplies), and within it class 6531, which
covers pharmaceutical medications subject to additional ceiling
regulation by the C\^amara de Regula\c{c}\~ao do Mercado de
Medicamentos (CMED). The three-letter acronyms matter: CAUFESP
defines who can bid, CADMAT defines what is being bid on, and
CMED caps what some of it can cost.

\paragraph{Enforcement of the SME-only rule.}
What LC~147 converted from an option into a default, TCE-SP is
what turns the default into behavior. Enforcement in BEC operates
through \emph{opt-out costs}: a purchasing unit that wants to
run a general (non-restricted) tender for an eligible item must
file a formal justification with the Tribunal de Contas do
Estado---the state audit court, and a body with real teeth. The
filing is reviewed, and a unit found to have circumvented the
SME-only rule without adequate basis can face administrative
sanctions, reporting penalties, and loss of procurement autonomy.
Buyers notice. State-level adherence to SME-only tendering for
eligible items rose from roughly 13\% before LC~147 to close to
70\% within three years, and by 2017 an un-justified general
tender on a sub-threshold item had become unusual enough to draw
scrutiny on its own. Buyers keep their files clean by using the
rule, not by arguing with it.

The R\$80{,}000 value ceiling is only half the trigger. Art.~49
of LC~123 adds a second condition: at least three eligible SMEs
must be ``potentially available'' for the item. Fewer than three,
and the auction falls back to the general rule. This
\emph{three-SMEs precondition} creates a natural item-level fault
line that is partly endogenous to the item's market. It also
becomes the lever of the 2014--2018 exception described next.

\paragraph{The Group-65 exception (2014--2018).}
One product group did not look like the others. In the months
after LC~147 took effect in 2014, São Paulo's state government
and TCE-SP together argued that medical and hospital supplies
belonged in a different regulatory category---``strategic goods''
that warranted open competition on public-health grounds---and
the audit court issued a resolution carving Group 65 out of the
SME-only default. The argument rested on three supporting points.
First, medical procurement has unusual downstream stakes: the
difference between an adequately stocked hospital and a
stock-out is not the same as the difference between two office
supplies, and rigorous price discovery mattered. Second, ANVISA's
pharmaceutical registration system narrows the qualified-supplier
pool for many medical items below the three-SMEs threshold of
LC~123, which on this reading would have made the restriction
inapplicable to large shares of the group anyway. Third, CMED's
regulated list prices for class-6531 medications already impose
a price ceiling that constrains buyer welfare, making the
additional SME restriction redundant for the affected items.
The three arguments together left Group 65 uniquely outside
the SME-only regime from August 2014 through February 2018---the
sole major product group so exempted during the period, and,
on the data, a regime-free control that the rest of the
catalogue did not have.

\paragraph{The March 2018 reversal.}
The 2014 carve-out was a compromise between the buying side
(SEFAZ/SP) and the audit side (TCE-SP). The 2018 reversal came
from a third actor---the \emph{Procuradoria-Geral do Estado}
(PGE-SP), the state's constitutional and administrative legal
counsel---which had not been at the original table. In November
2017, in the course of a routine internal legal review, PGE-SP
issued a formal opinion (parecer) concluding that the Group-65
exception violated \emph{isonomia}, the equal-treatment
principle codified in Art.~5 of the Federal Constitution. The
opinion did not engage with procurement outcomes, price
trajectories, or supplier complaints; it argued from first
principles that procedural symmetry across product groups
required a substantive distinguishing feature, and found that
the 2014 rationale did not clear that bar. ANVISA's registration
system, the opinion held, operates at the item level and is
already addressed by the three-SMEs precondition built into
LC~123; CMED's price controls operate as a ceiling rather than a
competitive floor, leaving the question of allocative efficiency
untouched. The remedy, PGE-SP directed, was to extend the
SME-only default to Group 65 on the same footing as every other
group. The parecer was transmitted to BEC in early 2018 and took
effect in March of that year. From the bidders' perspective, a
category that had been open for more than three years became,
overnight, another line in the SME-only catalogue.

\paragraph{Why this is a clean natural experiment.}
Three features of the March 2018 reversal support the identifying
assumption that the cut is exogenous to the procurement process
itself. First, the \emph{trigger} was a legal reinterpretation of
a constitutional principle, not a response to market outcomes:
no procurement-data audit, political event, or supplier complaint
preceded the opinion, and the parecer cites none. Second, the
\emph{timing} was bureaucratic. PGE-SP's opinion was issued on
its own internal schedule, transmitted to BEC in early 2018, and
took effect in March 2018 without the sort of public consultation
or draft-rule window that would have let bidders re-price in
advance. Third, the \emph{boundaries} of the experiment are
legally clean: the affected set is a single well-defined product
group (Group 65, meeting the three-SMEs precondition), and the
unaffected controls are the other 76 product groups that had
been operating under the SME-only default since 2014. An
18-month window straddling the break yields a crisp
treated-vs-never-treated comparison.

These features motivate the identifying strategy that follows in
Section~\ref{sec:identification}. The reversal should leave the
non-SME cost primitive alone; only the set of bidders that
participate changes, not who any given firm is. That invariance
is the core identifying assumption of this paper, and---crucially
for what follows---it is testable directly against the data.

%% ============================================================================
\section{Reduced-Form Context}
\label{sec:rf}

Before turning to primitives, two reduced-form facts motivate the
structural exercise: the March~2018 regime switch causes a 10--11\%
price increase, and a Gelbach mediation decomposition locates
\emph{95\%} of it outside any observable channel. This section
establishes both.

\paragraph{Design and estimate.}
March 2018 places Group~65 (the treated cohort) under the
SME-only default that had already applied to 76 other product
groups since 2014. Identification proceeds via a staggered-DiD
reversal: Group~65 switches from ``open'' into ``SME-only'' while
the 76 comparison groups remain SME-only throughout an 18-month
window. The \citet{goodmanbacon2021} decomposition places 100\%
of the two-way-fixed-effects weight on the clean
treated-vs-never-treated 2$\times$2 comparison---there is nothing
to re-aggregate. The main regression is
\begin{equation}
y_{ipt} \;=\; \beta\,(\mathit{g65}_i \times \mathit{Pre}_t)
\;+\; \zeta_i \;+\; \mathbf{X}'_{ipt}\gamma \;+\; \varepsilon_{ipt},
\label{eq:did-main}
\end{equation}
with item fixed effects $\zeta_i$, controls $\mathbf{X}$ including
a sealed-bid indicator and log quantity, and standard errors
clustered at the item level. Table~\ref{tab:rf-headline} reports
$\hat\beta$ across four outcomes and three windows.

\begin{table}[htbp]
\centering
\caption{Reduced-form headline DiD estimates (Group 65 vs 76 controls)}
\label{tab:rf-headline}
\begin{threeparttable}
\footnotesize
\setlength{\tabcolsep}{4pt}
\begin{tabular}{lrrrrrr}
\toprule
                  & \multicolumn{2}{c}{6-month} & \multicolumn{2}{c}{12-month} & \multicolumn{2}{c}{18-month} \\
                  & Baseline & + PBU FE & Baseline & + PBU FE & Baseline & + PBU FE \\
\midrule
Log winning price & $-0.142$ & $-0.148$ & $-0.108$ & $-0.108$ & $-0.109$ & $-0.113$ \\
                  & (0.015)  & (0.014)  & (0.013)  & (0.013)  & (0.012)  & (0.012)  \\
Log firm count    & $+0.211$ & $+0.220$ & $+0.128$ & $+0.138$ & $+0.098$ & $+0.107$ \\
                  & (0.009)  & (0.009)  & (0.007)  & (0.008)  & (0.007)  & (0.007)  \\
Log valid bids    & $+0.203$ & $+0.215$ & $+0.103$ & $+0.114$ & $+0.069$ & $+0.080$ \\
                  & (0.013)  & (0.014)  & (0.010)  & (0.011)  & (0.010)  & (0.011)  \\
Winner-PBU dist.\ (km) & $+6.64$ & $+1.94$ & $+13.78$ & $+5.70$ & $+14.25$ & $+6.15$ \\
                  & (3.32)  & (1.98)  & (3.02)  & (2.38)  & (3.08)  & (2.22)  \\
\bottomrule
\end{tabular}
\begin{tablenotes}
\footnotesize
\item \textit{Notes:} Each cell is $\hat\beta$ (SE) from
Equation~\eqref{eq:did-main} for the outcome-window-specification
combination. Item fixed effects included throughout. ``Baseline''
columns include sealed-bid and log-quantity controls; ``+ PBU FE''
columns add buyer-unit fixed effects. SE clustered at the item
level. Sample sizes vary by outcome and window
($\sim$250k--700k item-level observations). All price and
firm-count coefficients are significant at $p < 0.01$. Own
calculation on BEC bid-level data; replication scripts in the
supplementary materials.
\end{tablenotes}
\end{threeparttable}
\end{table}

Four patterns emerge from the 18-month baseline.
\textbf{Winning prices rise 10--11\%} (using $\exp(\hat\beta) - 1$;
sharper in shorter windows). \textbf{Firm count falls 10--11\%}
despite a policy mandating SME participation---non-SMEs exit
faster than SMEs enter. \textbf{Valid bids fall only 6--8\%},
consistent with each remaining SME submitting more bids per
auction than the non-SMEs they replaced. \textbf{Winner-to-PBU
distance falls 5--14 km}, consistent with SMEs being more
geographically local. The price coefficient is stable across the
Callaway-Sant'Anna, Sun-Abraham, and Goodman-Bacon
heterogeneity-robust estimators ($-0.108$ to $-0.114$);
cluster-robust to item-group, PBU, and two-way clustering; wild
cluster bootstrap delivers $p < 0.001$; a synthetic-control
comparison confirms the post-treatment gap; a formal spillover
test bounds implied bias at $+0.0005$ log-points; a
\citet{rambachan2023} HonestDiD diagnostic rules out plausible
parallel-trends violations.

\paragraph{Gelbach (2016) decomposition: the puzzle.}
The natural first question is whether the price jump is
compositional---fewer bidders, different types. A
\citet{gelbach2016} decomposition of the price regression adds
log firm count and a winner-SME indicator as mediators, asking
how much of $\hat\beta$ they jointly explain. The answer is
striking: \emph{almost none of it}. The two mediators move in
opposite directions---log firm count pulls prices down by
$+181\%$ of the headline, winner-SME composition pulls them up
by $-81\%$---and what looks like a large decomposition is a near
cancellation. The net mediated share is about 5\%; the
residual---the share of the price effect that observable entry
and observable composition cannot explain---is about
\textbf{95\%}. That residual is the puzzle the rest of the paper
addresses.

\paragraph{From reduced form to structural.}
The Gelbach cancellation rules out easy compositional stories.
What remains must be one (or more) of: (i) within-auction
bidding behavior that shifts with the candidate pool, (ii)
unobserved compositional shifts, or (iii) selection of firms
into bidding with different cost primitives. The reduced form
cannot separate these; the CPV inversion can. The rest of the
paper pulls the 95\% residual apart into a within-auction
intensive margin ($\sim$66\%), a cost-primitive entry margin
($\sim$34\%), a documented participation response that amplifies
the policy's reach without driving the price effect, and a
welfare calculation in both R\$ and US\$. Extended reduced-form
analysis---placebos, spillover bounds, policy-design variants---
is developed in a parallel paper~\citep{genicolo2026a}; the
present paper is self-contained on all quantities it uses. The
reduced form tells us where to look; the structural model tells
us what we are looking at.

%% ============================================================================
\section{A Model of Asymmetric FPSB Procurement with Type-Specific Costs}
\label{sec:model}

\paragraph{Primitives.}
The environment is a first-price sealed-bid reverse (procurement)
auction for a single indivisible good, with two bidder types that
match the institutional distinction the policy draws: type-A
(\emph{SME}) and type-B (\emph{non-SME}). In auction $t$, there are
$n_A^t \in \{0, 1, 2, \ldots\}$ SMEs and $n_B^t \in \{0, 1, 2,
\ldots\}$ non-SMEs. Type-$k$ bidders draw private costs $c \sim F_k$
i.i.d.\ of each other and of auction characteristics. The
distributions $F_A$ and $F_B$ are continuous with densities $f_A$,
$f_B$ on $[\underline{c}, \overline{c}]$. The two type-specific
distributions are the objects the paper is ultimately trying to
recover from the data, and the separation between them is what
allows the SME-only policy to be a policy rather than a
relabeling.

Each bidder observes only her own cost and the auction's type
composition $(n_A^t, n_B^t)$. Each submits a sealed bid
$b \in [0, \overline{c}]$; the lowest bid wins; ties are broken by a
fair lottery. The payoff to a type-$k$ bidder with cost $c$ who bids
$b$ is $(b - c) \cdot \mathbf{1}\{b \text{ wins}\}$.

\paragraph{Equilibrium.}
In a symmetric-within-type Bayes-Nash equilibrium, each type-$k$ bidder
uses the same bidding function $\beta_k: [\underline{c}, \overline{c}]
\to \mathbb{R}_+$, which is strictly increasing and continuously
differentiable on the interior of its support. Let $G_k(b) = \Pr(b_k \le b)$
denote the induced bid distribution and $g_k(b)$ its density. These
are observables from auction data.

A type-$k$ bidder with cost $c$ maximizes
\[
\pi_k(b, c; n_A^{-k}, n_B^{-k})
= (b - c) \cdot
  \bigl[1 - G_A(b)\bigr]^{n_A^*}
  \bigl[1 - G_B(b)\bigr]^{n_B^*},
\]
where $n_A^* = n_A - \mathbf{1}\{k = A\}$, $n_B^* = n_B - \mathbf{1}\{k = B\}$
denote the \emph{opponent} counts of each type.

The first-order condition is
\begin{equation}
b - c \;=\; \frac{1}
  {\;n_A^* \, \dfrac{g_A(b)}{1-G_A(b)}
   + n_B^* \, \dfrac{g_B(b)}{1-G_B(b)}\;},
\label{eq:cpv-foc}
\end{equation}
and the inverse equilibrium bid function, central to identification, is
\begin{equation}
c \;=\; b \;-\; \frac{1}
  {\;n_A^* \, \dfrac{g_A(b)}{1-G_A(b)}
   + n_B^* \, \dfrac{g_B(b)}{1-G_B(b)}\;}.
\label{eq:cpv-inverse}
\end{equation}
Equation~\eqref{eq:cpv-inverse} is the asymmetric extension of
\citet{gpv2000} developed by \citet{campo2003}, and it is the
single piece of the model that does the estimation work. The
symmetric-within-type special case ($G_A = G_B = G$, $g_A = g_B =
g$) collapses to the familiar $c = b - (1 - G(b)) / ((N-1) g(b))$,
the standard procurement GPV. Crucially, once type-specific bid
distributions are estimated from the data, the inverse map is
deterministic: there are no parameters to search over, no auxiliary
distributions to assume. Whatever asymmetry exists between SME and
non-SME costs shows up directly in the pseudo-costs. Appendix~F
derives Equation~\eqref{eq:cpv-inverse} from the bidder's
first-order condition for readers who want the self-contained
derivation.

\paragraph{Ref-price normalization.}
A practical obstacle is that reference prices $P_{\text{ref}}^t$
span two orders of magnitude across Group-65 auctions, from small
consumables to large-volume medication purchases. Running CPV on
raw bids would mix scales that are not, in any economic sense,
comparable. Following the bid-preference literature
\citep[e.g.,][]{athey2013}, I normalize each bid by its
auction-specific reference price, $\tilde{b}^t = b^t /
P_{\text{ref}}^t$, placing all normalized bids on a common
$[0, 2]$ scale. Equation~\eqref{eq:cpv-inverse} then applies to
$\tilde{b}$, with $\tilde c = c / P_{\text{ref}}$ on the same
scale. The normalization amounts to an assumption that the cost
distribution scales linearly with the buyer's reference price---a
strong but standard assumption, and one that the invariance test
in Section~\ref{sec:estimation} provides an indirect check on.
Level results without normalization are reported in the appendix.

%% ============================================================================
\section{Identification}
\label{sec:identification}

\paragraph{Point identification from bid data.}
The CPV inversion is, at root, a one-line recipe. The model
implies that $G_A$, $G_B$, $g_A$, $g_B$ are directly observable
from type-separated bid data, and that is what the bid-level
records on BEC provide. Given these four objects and an auction's
type composition $(n_A^t, n_B^t)$, Equation~\eqref{eq:cpv-inverse}
returns a pseudo-cost $\hat c_{i,t}$ for every bidder $i$ in every
auction $t$. Repeating the operation across all bidder-auctions
delivers type-specific empirical cost distributions $\hat F_A$
and $\hat F_B$. Nothing more is needed: there is no nonlinear
search, no auxiliary distributional assumption, no second
estimation stage.

Two conditions have to hold for the recipe to be meaningful.
First, asymmetric bidders must actually appear in the same
auctions, so that both $G_k$ are estimable on a non-degenerate
sample. Convite Group 65 satisfies this condition comfortably in
the Pre period---40\% of auctions mix SMEs and non-SMEs. Second,
bid functions must be monotonic in cost within type, which the
model guarantees and which the data confirm (the pseudo-costs
satisfy $\hat c \leq b$ on every observation in the estimation
sample).

\paragraph{Identification check: primitive invariance.}
The policy gives us a second handle that most structural auction
papers do not have: a testable over-identifying restriction. In
the model, $F_A$ and $F_B$ are \emph{primitives}---they describe
the cost process before any equilibrium considerations enter.
The March 2018 reversal changes \emph{which} non-SMEs are eligible
to bid in Group 65 (none, post-policy), but it cannot reach inside
any given firm and change its cost draw. Therefore,
\[
F_B^{\text{Pre}}(c) \;=\; F_B^{\text{Post}}(c)
\quad \text{for all } c,
\]
and a Kolmogorov-Smirnov test of exactly this equality is the
paper's primary identification check. No analogous condition holds
for the SME distribution: the policy changes the return to entering
Group 65 for small firms, and we expect $F_A^{\text{Post}}$ to
shift relative to $F_A^{\text{Pre}}$ precisely because marginal
SMEs enter (or exit) under the restriction. Selection into the
SME pool is part of the entry channel we want to measure. The
over-identifying restriction is therefore one-sided by design:
$F_B$ invariant (testable), $F_A$ allowed to move.

\paragraph{Robustness to mis-specified equilibrium.}
The price of using an equilibrium condition to recover primitives
is that the recovery is only as good as the equilibrium. If
bidders are risk-averse rather than risk-neutral, or if the
auction has a common-value component, the pseudo-costs returned
by Equation~\eqref{eq:cpv-inverse} are not consistent estimators
of the underlying cost distribution. The invariance test on
$F_B$ is therefore a joint test of the model \emph{and} of
primitive invariance: passing it is evidence for both. Two useful
features soften this concern in practice. First, if the
mis-specification is common to both bidder types, the invariance
condition still identifies a weaker but coherent hypothesis about
the unchanged non-SME pool. Second, the risk-aversion and
common-value concerns are addressed directly in
Section~\ref{sec:robustness} using a CARA CGPV calibration and
the \citet{hendricks2003} diagnostic, respectively. Results in
Section~\ref{sec:estimation} suggest the model is reasonably
specified: invariance holds tightly for non-SMEs (the type whose
pool is supposed to be unchanged), and fails in exactly the
direction the entry channel predicts for SMEs.

%% ============================================================================
\section{Data}
\label{sec:data}

\paragraph{BEC bid-level records.}
Because every São Paulo state purchase passes through a single
electronic platform, the data requirements of a structural auction
paper sit unusually well with the data actually available. I use
the complete bid-level trail on BEC for 2009--2019: 39.96 million
individual bids across 601{,}229 \emph{Preg\~ao eletr\^onico}
auctions, 612{,}120 \emph{Convite} sealed-bid auctions, and
140{,}679 \emph{Dispensa eletr\^onica} auctions. Each bid record
is a row in the platform's transaction log and carries the
bidder's tax ID (CNPJ), bid value, auction identifier, item
catalog code, administrative reference price, winner flag, and
timestamps. The same platform that enforces the procurement rules
also records every keystroke in a structured form---an
uncommonly clean dataset for a developing-country setting.

\paragraph{Convite modality focus.}
The structural estimation uses the \emph{Convite} modality, and
the reason is institutional as much as technical: Convite is the
only modality in BEC that operationally implements a first-price
sealed-bid auction. Three features of its design line up with
what FPSB identification needs.

\noindent(1) \textbf{One sealed envelope per firm.} Convite
requires each invited bidder to submit exactly one envelope, and
iterative bidding is disallowed by statute. In the raw data,
99.9\% of firm-auction pairs have exactly one bid; the remaining
0.1\% are administrative re-submissions in re-opened auctions
(Art.~48~\S3, Lei~8.666/93) and are trimmed at the firm-auction
level. The observed bid is therefore the firm's equilibrium bid,
not a step in an iterative process.

\noindent(2) \textbf{Mixed-type participation.} In the Pre period,
40\% of Convite Group-65 auctions have at least one SME and at
least one non-SME bidding in the same auction. This is exactly
the identifying variation CPV requires: if every auction were
single-type, we could only estimate the two $G_k$'s in isolation,
without the cross-type competition that gives the asymmetric
model its bite.

\noindent(3) \textbf{Observed reference price.} The administrative
reference price $P_{\text{ref}}$ is set by the purchasing unit
before the auction opens and is recorded on 98.7\% of Group-65
Convite auctions. This is what makes the ref-price normalization
from Section~\ref{sec:model} operational in practice: every
auction carries a scale-free anchor that the platform itself
records.

\paragraph{Sample restrictions.}
I restrict the structural sample to:
\begin{itemize}
  \item Group~65 (medical supplies, the treated product group);
  \item the 18-month analytical window Sep 2016 -- Aug 2019, centered on
        the March 2018 regime break;
  \item auctions with $N \ge 2$ bidders (GPV estimator undefined at $N=1$);
  \item normalized bids $\tilde b \in (0.005, 2)$ (discard R\$0.01
        placeholder bids and bids implausibly above reference).
\end{itemize}
The pipeline yields two nested samples. The \emph{CPV-eligible}
sample---Convite Group-65 auctions with $N \geq 2$ bidders within
the 18-month window---contains 33{,}706 auctions and 120{,}463
bidder-auctions (mean $N = 3.57$). Applying the CPV inversion
with valid reference price and normalized bids in $(0.005, 2)$,
and retaining pseudo-costs that satisfy the clean-filter
$\hat c / P_{\text{ref}} \in (0.001, 1.5)$, delivers the
\emph{structural estimation sample} of 24{,}346 bidder-auctions
in 24{,}336 distinct auctions. The CPV density estimation
pools bids across auctions within period-$N$ strata, so the
1-clean-pseudo-cost-per-auction post-filter retention rate does
not constrain identification---what matters is the density over
the full pooled bid distribution, which is estimated from the
pre-inversion population. Table~\ref{tab:sample} reports the
breakdown for both samples.

\begin{table}[htbp]
\centering
\caption{Structural estimation sample (Convite, Group~65)}
\label{tab:sample}
\begin{threeparttable}
\small
\setlength{\tabcolsep}{4pt}
\begin{tabular}{lrrr}
\toprule
Cell & Bidder-auctions & Auctions & Mean $N$ \\
\midrule
\textit{CPV-eligible sample} ($N \geq 2$)    & 120{,}463 & 33{,}706 & 3.57 \\
\hspace{1em} Pre-period                      & 51{,}068  & 14{,}411 & 3.54 \\
\hspace{1em} Post-period                     & 69{,}395  & 19{,}295 & 3.60 \\
\midrule
\textit{Structural sample} (post-inversion)  & 24{,}346  & 24{,}336 & 3.75 \\
\hspace{1em} Pre-period                      & 10{,}704  & 10{,}699 & 3.71 \\
\hspace{2em} SME bidder-auctions             & 3{,}624   & --       & -- \\
\hspace{2em} Non-SME bidder-auctions         & 7{,}080   & --       & -- \\
\hspace{1em} Post-period                     & 13{,}642  & 13{,}637 & 3.79 \\
\hspace{2em} SME bidder-auctions             & 6{,}786   & --       & -- \\
\hspace{2em} Non-SME bidder-auctions         & 6{,}856   & --       & -- \\
\midrule
\multicolumn{4}{l}{\textit{CPV-eligible sample, by auction size:}} \\
\hspace{1em} Auctions with $N = 2$           & -- & 13{,}010 & 2.00 \\
\hspace{1em} Auctions with $N = 3$           & -- & 8{,}360  & 3.00 \\
\hspace{1em} Auctions with $N = 4$           & -- & 4{,}762  & 4.00 \\
\hspace{1em} Auctions with $N \geq 5$        & -- & 7{,}574  & 6.64 \\
\bottomrule
\end{tabular}
\begin{tablenotes}
\footnotesize
\item \textit{Notes:} Convite (sealed-bid) auctions of Group~65
items in the 18-month window Sep 2016 -- Aug 2019. The
``CPV-eligible'' sample imposes $N \ge 2$ distinct bidders per
auction. The ``structural estimation'' sample further requires
valid reference price and normalized bid in $(0.005, 2)$ at the
bid level, and retains post-inversion pseudo-costs satisfying
$\hat c / P_{\text{ref}} \in (0.001, 1.5)$. ``Mean $N$ per
auction'' is the auction-level average bidder count (the same
$N$ that enters the CPV opponent-count in
Equation~\ref{eq:cpv-inverse}). SME classification is the
firm-registry proxy described in Appendix~A:
$\mathbf{1}\{$\texttt{porte\_empresa} $\in \{01, 03\}$ or
\texttt{fornec\_enquad} $\in \{1, 2\}\}$, which captures ME
(micro-empresa) and EPP (empresa de pequeno porte)
classifications.
\end{tablenotes}
\end{threeparttable}
\end{table}

%% ============================================================================
\section{Estimation}
\label{sec:estimation}

\paragraph{Procedure.}
The mechanics are deliberately plain. For each stratum
$(\text{period}, N\text{-bin})$, I estimate the type-specific bid
distributions $\hat G_A(\tilde b), \hat G_B(\tilde b)$ and
densities $\hat g_A(\tilde b), \hat g_B(\tilde b)$ by Gaussian
kernel density estimation on the pooled normalized bids of each
type, using Silverman's rule-of-thumb bandwidth
$h_k = 1.06 \cdot \hat\sigma_k \cdot n_k^{-1/5}$. CDFs and
densities are evaluated at each observed bid by linear
interpolation of the kernel estimates, and $(1-G)$ is floored at
$10^{-6}$ to keep the denominator in Equation~\eqref{eq:cpv-inverse}
well-defined near the upper tail of the bid distribution. None
of the bandwidth or floor choices is material: appendix
robustness replaces Silverman with a plug-in bandwidth and
tightens the floor by two orders of magnitude, and the headline
numbers move by less than a percentage point.

Equation~\eqref{eq:cpv-inverse} is then applied firm-by-firm,
auction-by-auction, using the auction-specific opponent counts
$(n_A^*, n_B^*)$ to recover a pseudo-cost $\hat c_{i,t}$ for
every firm-auction pair. The data are split into eight strata,
$(\text{Pre}, \text{Post}) \times \{N=2, N=3, N=4, N \ge 5\}$, so
that each inversion is applied within a roughly homogeneous
competitive environment. A small share of pseudo-costs fall
below zero as an artifact of boundary bias in the kernel
estimator; I flag them but keep them in the primary results.
Restricting to $\hat c \in (0, 1.5)$ retains 77.7\% of the
estimation sample, and all invariance and decomposition results
in what follows are computed on that clean subset.

\paragraph{Equilibrium-condition check.}
Before looking at the costs, it is worth verifying that they
obey the sign condition the model requires. A correctly specified
equilibrium delivers $\hat c \leq b$ at every firm-auction pair:
bidders cannot earn negative markups. Across all eight strata,
zero percent of pseudo-costs violate this condition. The CPV
inversion is numerically stable on this dataset, and the
estimator is not fighting the data.

\paragraph{Cost distributions and markups.}
With the pseudo-costs in hand, the first question is what the
type-specific distributions actually look like. Table
\ref{tab:cost-summary} summarizes.

\begin{table}[htbp]
\centering
\caption{CPV-recovered cost distributions (Convite Group~65; $c/P_{\text{ref}}$ units)}
\label{tab:cost-summary}
\begin{threeparttable}
\small
\begin{tabular}{llrrrrr}
\toprule
Period & Type & $n$ & Mean $c/P_{\text{ref}}$ & Median &
  Median markup & Markup \% \\
\midrule
Pre   & Non-SME   & 5{,}651 & 0.523 & 0.515 & 0.119 & 18.9\% \\
Pre   & SME       & 2{,}604 & 0.520 & 0.513 & 0.138 & 21.1\% \\
Post  & Non-SME   & 5{,}452 & 0.524 & 0.532 & 0.099 & 15.7\% \\
Post  & SME       & 5{,}200 & 0.560 & 0.572 & 0.131 & 18.6\% \\
\bottomrule
\end{tabular}
\begin{tablenotes}
\footnotesize
\item \textit{Notes:} Pseudo-costs recovered via CPV inversion
(Equation~\ref{eq:cpv-inverse}) with Gaussian-kernel density
estimation and Silverman bandwidth. Sample restricted to
$\hat c/P_{\text{ref}} \in (0.001, 1.5)$ to discard boundary-bias
artifacts. Markup $= b - c$; Markup\% $= (b-c)/b$.
\end{tablenotes}
\end{threeparttable}
\end{table}

Three facts stand out, and together they tell the story that
organizes the rest of the paper.

\noindent(i) \emph{Pre-period, SME and non-SME cost distributions
are virtually identical}. Mean $0.520$ versus $0.523$ of reference
price; medians $0.513$ versus $0.515$. Figure~\ref{fig:cpv-pre}
plots both empirical CDFs, and they overlap across the full
support. This is not a neutral finding: it means that under the
Open regime, the two bidder types were effectively drawing from
the same cost pool, even though regulators treat them as
categorically different.

\noindent(ii) \emph{Post-period, the non-SME distribution is
essentially unchanged} (mean $0.524$), while the SME distribution
shifts upward by 7.7\% relative to its Pre mean (from $0.520$ to
$0.560$; median $0.513$ to $0.572$).
One type's pool holds steady across the break; the other's moves
in exactly the direction a binding entry constraint would
predict.

\noindent(iii) \emph{Markups decrease monotonically in $N$},
as the theory requires---the check against the CPV output
behaving pathologically (Table~\ref{tab:markup-by-n},
Appendix~B).

\begin{figure}[htbp]
\centering
\includegraphics[width=0.72\textwidth]{fig_v2_cpv_costs_pre.pdf}
\caption{Under the open regime, SME and non-SME cost distributions
are virtually indistinguishable (Pre-period Convite G65). The two
empirical CDFs overlap across the full support, undermining the
standard justification for the SME-only set-aside as a level-the-
playing-field instrument: at the margin, the ``playing field'' was
already level.}
\label{fig:cpv-pre}
\end{figure}

\paragraph{Primitive-invariance test.}
The over-identifying restriction from Section~\ref{sec:identification}
can be evaluated directly. Table~\ref{tab:invariance} reports
Kolmogorov-Smirnov tests of the hypothesis $F_k^{\text{Pre}} =
F_k^{\text{Post}}$ for each type---the central check that
separates the model's identifying content from a tautology.

\begin{table}[htbp]
\centering
\caption{Primitive-invariance test ($F_k^{\text{Pre}} = F_k^{\text{Post}}$)}
\label{tab:invariance}
\begin{threeparttable}
\small
\begin{tabular}{lrrrrrr}
\toprule
Type & $n^{\text{Pre}}$ & $n^{\text{Post}}$ & Mean Pre & Mean Post
  & KS $D$ & KS $p$ \\
\midrule
Non-SME (should be invariant) & 5{,}651 & 5{,}452 & 0.523 & 0.524 & 0.030 & 0.012 \\
SME (should shift)            & 2{,}604 & 5{,}200 & 0.520 & 0.560 & 0.081 & $<10^{-10}$ \\
\bottomrule
\end{tabular}
\begin{tablenotes}
\footnotesize
\item \textit{Notes:} Invariance test under the null that the cost
distribution is primitive and therefore unaffected by the regime switch.
Non-SME passes the test economically (shift of $0.001$ is below any
reasonable bias threshold); SME fails, consistent with the entry of
marginal high-cost SMEs under the binding set-aside regime post-March 2018.
\end{tablenotes}
\end{threeparttable}
\end{table}

The numbers are cleaner than one usually has any right to expect.
The non-SME distribution is economically invariant across the
March 2018 break: the mean shifts by 0.001 of reference price
(the KS statistic marginally rejects formal equality at
$p = 0.012$, but this is a consequence of the 11{,}000-observation
sample size, not of any substantive shift). The SME distribution,
by contrast, moves by 0.039 ref-price units upward, with
$p < 10^{-10}$. Read together, this is the sharpest form the data
can take: the model's identifying restriction holds where it is
supposed to hold, and fails in exactly the direction---and on
exactly the type---that the entry channel predicts.

Figure~\ref{fig:cpv-stab} shows the two panels visually.

\begin{figure}[htbp]
\centering
\includegraphics[width=0.85\textwidth]{fig_v2_cpv_stability.pdf}
\caption{Primitive-invariance check: non-SME cost distribution is
   invariant across the March~2018 break (left panel); SME cost
   distribution shifts upward, consistent with marginal entry
   (right panel).}
\label{fig:cpv-stab}
\end{figure}

%% ============================================================================
\section{Structural Decomposition of the Price Effect}
\label{sec:decomposition}

\paragraph{Counterfactual scenarios.}
With type-specific primitives in hand, the decomposition becomes
a matter of asking the model three slightly different questions
about the same population of auctions. For each Pre-period
auction with observed composition $(n_A^t, n_B^t)$ and reference
price $P_{\text{ref}}^t$, I simulate three scenarios:

\begin{description}
\setlength\itemsep{0.5em}
  \item[Scenario S1 (Open, Pre baseline)]
    all $n_A^t + n_B^t$ bidders participate with their Pre-period cost
    distributions, with expected winning bid
    \[
    E[b_{\min} | \text{S1}] = \int \bar G_A^{\text{Pre}}(\tilde b)^{n_A}
    \bar G_B^{\text{Pre}}(\tilde b)^{n_B} \, d\tilde b,
    \]
    where $\bar G_k = 1 - G_k$ is the empirical survival function.
  \item[Scenario S2 (SME-only, Pre pool)]
    only the $n_A^t$ SMEs from the Pre pool bid, with
    \[
    E[b_{\min} | \text{S2}] = \int \bar G_A^{\text{Pre}}(\tilde b)^{n_A}
    \, d\tilde b.
    \]
  \item[Scenario S3 (SME-only, Post pool)]
    $n_A^t$ SMEs bid, drawn from the Post-period SME cost
    distribution (which includes marginal entrants), with
    \[
    E[b_{\min} | \text{S3}] = \int \bar G_A^{\text{Post}}(\tilde b)^{n_A}
    \, d\tilde b.
    \]
\end{description}

The decomposition follows directly: the \emph{intensive margin}
is $\Delta^{\text{int}} = E[b|\text{S2}] - E[b|\text{S1}]$, the
effect of shutting non-SMEs out at a fixed Pre SME pool; the
\emph{entry margin} is $\Delta^{\text{ent}} = E[b|\text{S3}] -
E[b|\text{S2}]$, the additional effect of letting the SME pool
re-compose toward the post-policy mix; and the two sum to the
total structural effect $\Delta^{\text{tot}} = E[b|\text{S3}] -
E[b|\text{S1}] = \Delta^{\text{int}} + \Delta^{\text{ent}}$. S1,
S2, S3 are integrals of observed type-specific bid distributions
rather than re-simulated equilibrium bid functions, which means
$\Delta^{\text{int}}$ is a \emph{conservative} measurement of the
within-auction channel: a fully re-solved equilibrium with Pre
SME cost primitives but an all-SME competitive environment would
predict SMEs bidding \emph{softer} than their observed Pre bids,
because they face fewer low-cost opponents. The intensive share
reported here is therefore a lower bound on the true
bid-function response; the 66/34 split understates, not
overstates, the intensive margin's dominance.

\paragraph{Results.}
Table~\ref{tab:decomp} reports the decomposition by auction-size bin.

\begin{table}[htbp]
\centering
\caption{Structural decomposition of the SME-only price effect}
\label{tab:decomp}
\begin{threeparttable}
\small
% Auto-generated by 36_counterfactuals.R
\begin{tabular}{lrrrrrrrr}
\toprule
N-bin & Auctions & $E[b^{S1}]$ & $E[b^{S2}]$ & $E[b^{S3}]$ & Intensive & Entry & Total & Intens.\% \\
\midrule
N=2 & 5,526 & 0.465 & 0.563 & 0.597 & 0.099 & 0.034 & 0.133 & 74\% \\
N=3 & 3,656 & 0.357 & 0.460 & 0.575 & 0.102 & 0.115 & 0.218 & 47\% \\
N=4 & 2,082 & 0.310 & 0.454 & 0.515 & 0.143 & 0.062 & 0.205 & 70\% \\
N>=5 & 3,147 & 0.273 & 0.391 & 0.438 & 0.118 & 0.046 & 0.165 & 72\% \\
\bottomrule
\end{tabular}

\begin{tablenotes}
\footnotesize
\item \textit{Notes:} All quantities in $c/P_{\text{ref}}$ units.
Intensive = $E[b|\text{S2}] - E[b|\text{S1}]$; Entry =
$E[b|\text{S3}] - E[b|\text{S2}]$; Total = $E[b|\text{S3}] -
E[b|\text{S1}]$. Intens.\% = share of total attributable to the
intensive margin. Simulation uses the observed distribution of
$(n_A^t, n_B^t)$ in Pre-period Convite G65 auctions within each
$N$-bin. Bootstrap standard errors on these shares are reported
in Appendix~G (Table~\ref{tab:bootstrap-se}).
\end{tablenotes}
\end{threeparttable}
\end{table}

Under this fixed-pool accounting, the intensive margin dominates.
It captures 47--74\% of the total structural price effect across
$N$-bins, with an (unweighted) average of about two-thirds: the
within-auction bidding response explains roughly twice as much of
the policy's price effect as the pool-composition response does.
The entry channel peaks at $N=3$ (53\%) and bottoms out at $N=2$
(26\%), with no clean monotone pattern---the entry margin is not
uniformly larger in busier auctions or thinner ones, it follows
the local shape of the cost distributions. On an auction-weighted average,
intensive absorbs 66\% of the total structural price effect and
entry absorbs the remaining 34\%.

Figure~\ref{fig:decomp} plots the decomposition as stacked bars.

\begin{figure}[htbp]
\centering
\includegraphics[width=0.80\textwidth]{fig_v2_decomp.pdf}
\caption{Stacked decomposition of the SME-only price effect by
   auction-size bin. The intensive margin (dark grey) dominates;
   entry (light grey) matters most at $N=3$.}
\label{fig:decomp}
\end{figure}

\paragraph{Reconciliation with the Gelbach residual.}
It is now possible to go back to the 95\% residual from
Section~\ref{sec:rf} and say what was inside it. The reduced-form
Gelbach decomposition could not measure either of the two channels
the structural exercise isolates: it used \emph{observed} entry
(the number of bidders) and \emph{observed} winner composition
(whether the winner was an SME), but neither object is the same
as the cost-primitive entry that shows up in
$F_A^{\text{Pre}} \neq F_A^{\text{Post}}$, and neither captures
the within-auction bidding response the CPV inversion makes
visible. The structural decomposition breaks the 95\% residual
into 66\% within-auction bidding response and 34\% cost-primitive
entry. The dominant channel is the one no reduced-form mediator
could have seen---the bid-function adjustment that a firm makes
when it realizes it is no longer facing a non-SME competitor.
Participation also responds substantially, as
Section~\ref{sec:endogenous} documents, but the magnitude of the
participation response is a separate fact about the policy's
reach; it is not the mechanism through which the price effect
itself is generated.

%% ============================================================================
\section{The Participation Response}
\label{sec:endogenous}

The 66/34 decomposition of Section~\ref{sec:decomposition} is the
paper's structural headline. It answers a specific question: given
the bidders who actually showed up, how much of the 10--11\% price
effect comes from how they bid versus who they are? This section
documents a second, complementary fact about the policy---that
participation itself responds substantially to the regime
change---and shows that this response, while real, does not
overturn the within-auction primacy of the intensive margin. The
policy reorganizes who shows up, but the price movement is
primarily generated by how the remaining bidders bid.

\paragraph{Reduced-form entry regression.}
A simple descriptive exercise makes the scale of the entry
response visible before any structural model is imposed on it.
For each CADMAT item class, I define the \emph{potential bidder
pool} as the set of CNPJs that have bid on any item in the class
at some point in the 18-month window---a practical operationalization
of ``would plausibly consider this auction.'' For each auction,
the type-$k$ entry rate is the ratio of observed type-$k$ bidders
to the type-$k$ potential pool. Log entry rates are then
regressed on a Post indicator and a Post$\times$SME interaction
with class fixed effects:
\[
\log \text{entry\_rate}_{i,t,k}
 \;=\; \alpha_{c(i,t)}
    + \beta_1 \cdot \text{Post}_t
    + \beta_2 \cdot \text{Post}_t \times \mathbf{1}\{k = \text{SME}\}
    + \varepsilon_{i,t,k}
\]
with class fixed effects $\alpha_c$ and clustering at the auction
level (Table~\ref{tab:entry-reg}).

\begin{table}[htbp]
\centering
\caption{Entry-rate regression (Convite G65, 18-month window)}
\label{tab:entry-reg}
\small
\begin{tabular}{lrrrr}
\toprule
                     & Estimate & SE    & $t$     & $p$ \\
\midrule
Post                 & $-0.961$ & 0.023 & $-42.7$ & $<10^{-16}$ \\
Post $\times$ SME    & $+1.375$ & 0.037 & $+37.6$ & $<10^{-16}$ \\
Class fixed effects  &          &       &         & yes \\
Observations         & 66{,}720 &       &         & \\
Within $R^2$         & 0.031    &       &         & \\
\bottomrule
\end{tabular}
\end{table}

The magnitudes are large and tightly estimated. The non-SME
entry rate falls by $1 - \exp(-0.961) = 62\%$ under the Post
regime; the SME entry rate rises by $\exp(-0.961 + 1.375) -
\exp(-0.961) \approx +52\%$. Both responses are extremely
statistically significant ($|t| > 37$). On raw counts, mean
non-SMEs per auction drops from $2.05$ to $1.73$ ($-15.9\%$) and
mean SMEs rises from $1.49$ to $1.87$ ($+25.7\%$). The policy
does not merely filter who wins---it reorganizes who even bothers
to show up, on both sides of the distinction.

\paragraph{Implied entry cost under free entry.}
Under the \citet{atheyseira2011} free-entry zero-profit condition,
$E[\text{profit}_k | \text{entry}] = K^k$, the type-specific entry
cost is pinned down by observed markups and win probabilities.
Plugging in the CPV-recovered markups and observed winning
probabilities, the implied entry costs are R\$30.5 per auction
for non-SMEs and R\$33.3 for SMEs (US\$8.7 and US\$9.5 at R\$3.50/USD,
the 2016--2019 BCB closing-rate average). The near-equality is a useful
sanity check rather than a puzzle: the marginal cost of preparing
a Convite bid in BEC is essentially the same whichever type you
are---the CAUFESP registration, qualification certificates, and
bid-envelope preparation do not materially depend on firm size.
A free-entry model with type-invariant entry costs is therefore
not a heroic abstraction in this setting.

\paragraph{Does the participation response overturn the 66/34?}
A skeptical reading might ask whether the fixed-pool headline is
an artifact of treating participation as exogenous when it
visibly responds. A reduced-form calibration that credits the
observed non-SME exit and SME entry entirely to the entry channel
shifts the decomposition by construction: the within-auction
share drops from 66\% to roughly one-third and the participation
share rises correspondingly (details and the by-$N$ sensitivity
in Appendix~E, Table~\ref{tab:decomp-endogenous}). This is not,
however, evidence that the participation channel \emph{dominates}
the price effect. Three considerations anchor the headline in the
fixed-pool result.

The central point is that \emph{the participation shifts are
consistent with a large intensive-margin channel, not a substitute
for it}. The same re-selection that the calibration credits to
the entry channel also tightens the competitive environment for
the SMEs who remain in the pool: they face fewer bidders, and
they rationally bid softer in response. This softer bidding is
precisely what the fixed-pool CPV inversion measures on the
recovered Post-period bid distributions. In other words, the
CPV intensive margin already captures the bid-function response
that is \emph{caused by} the participation shifts; adding a
separate ``participation channel'' on top double-counts that
response. The calibration double-counts re-selection; the CPV
inversion does not.

Two supporting considerations reinforce this reading. First,
the calibration is a reduced-form attribution rather than a
primitive-level result. A rigorous allocation of price movement
to endogenous entry requires a nested fixed-point entry-bidding
equilibrium along the lines of \citet{atheyseira2011}, with
explicit distributions over entry costs and ex-ante signals.
The paper estimates entry costs as type-level scalars only.
Second, the primitive-invariance test that validates the CPV
estimator---$F_B^{\text{Pre}} = F_B^{\text{Post}}$---is a test
about the non-SME \emph{pool} as observed in Convite. It passes
tightly, confirming that CPV is recovering real cost primitives
rather than equilibrium-selection artifacts. The endogenous-entry
analogue---whether the marginal entrant's cost draw equals the
infra-marginal bidder's---is not tested directly and would require
an additional layer of modeling.

The participation response is therefore best read as
\emph{confirmatory context}: the policy reorganizes entry on top
of changing bidder behavior, and both effects raise the total
welfare cost the policy generates. The dominant price channel,
under the rigorous decomposition, is within-auction bidding.

\begin{figure}[htbp]
\centering
\includegraphics[width=0.75\textwidth]{fig_v2_entry_rates.pdf}
\caption{Entry rates by auction size, type, and period. Non-SME
rates fall markedly under the Post regime; SME rates rise. The
reorganization is large but, under the CPV decomposition,
contributes the minority of the price effect.}
\label{fig:entry-rates}
\end{figure}

%% ============================================================================
\section{Welfare Analysis}
\label{sec:welfare}

\subsection{Welfare Accounting Framework}
\label{sec:welfare-framework}

With type-specific cost primitives in hand, the welfare question
becomes arithmetically concrete: who pays for the restriction,
and how much. There are three candidates---the buyer (fiscal
cost), the remaining winners (rent change), and society
(deadweight loss)---and a simple accounting identity ties them
together. The identity is standard for procurement auctions, but
worth stating explicitly: with asymmetric bidder types, the signs
and magnitudes of its three pieces can pull in different
directions.

Consider a single auction with potential bidders of type
$k \in \{A, B\}$ and auction-level reference price $P_{\text{ref}}$.
Under any regime $R$, the winner is the bidder submitting the
lowest bid $b_{(1)}^R$; denote the winner's cost draw $c_{(1)}^R$. The
buyer's payment is $b_{(1)}^R$, the winner's profit is
$b_{(1)}^R - c_{(1)}^R$, and the total resource cost expended in
producing the good is $c_{(1)}^R$. Comparing two regimes $R, R'$ and
taking expectations over cost draws:
\begin{align}
\Delta_{R \to R'}\, \mathbb{E}[b_{(1)}] &= \Delta_{R \to R'}\, \mathbb{E}[c_{(1)}]
\;+\; \Delta_{R \to R'}\, \mathbb{E}[b_{(1)} - c_{(1)}],
\label{eq:welfare-identity}
\end{align}
The left-hand side is what the buyer pays more of. The right-hand
side splits that extra payment into two parts: a \emph{resource-cost
component}, which records whether the eligible pool can actually
produce the good at lower minimum cost, and a \emph{winner-rent
component}, which records whether the winning bidder gets to keep
more or less of the gap. The resource-cost piece is the deadweight
loss of the restriction. It is non-negative almost by definition:
taking bidders out of a set can only raise, never lower, the
minimum cost available within it. The rent piece is genuinely
ambiguous---when SMEs face only other SMEs their pool is more
symmetric and competitive, which tends to compress markups, so
winners can easily end up with \emph{less} profit per auction than
non-SMEs would have earned. Under the SME-only rule, then, the
buyer pays more while each winner keeps less, and the two
components do not cancel: their difference is the social loss.

Read in terms of standard welfare objects, identity
\eqref{eq:welfare-identity} has three entries:
\begin{itemize}
    \item \textbf{Fiscal cost} ($\Delta \mathbb{E}[b_{(1)}]$): what
      the buyer pays more, per auction, under the restriction. This
      is the same quantity the reduced-form regression in
      Section~\ref{sec:rf} estimates directly from observed winning
      bids.
    \item \textbf{Deadweight loss} ($\Delta \mathbb{E}[c_{(1)}]$):
      the excess real resources used to produce the good under the
      restriction---the part of the fiscal cost that is not a
      transfer between buyer and seller but a genuine loss.
    \item \textbf{Rent change} ($\Delta \mathbb{E}[b_{(1)} - c_{(1)}]$):
      the change in the winning firm's profit. Negative in our
      setting: SMEs competing against other SMEs face a tighter,
      more symmetric book than they did when non-SMEs were in the
      pool.
\end{itemize}
The structural primitives are what make this decomposition
operational. Without $F_c^A$ and $F_c^B$, identity
\eqref{eq:welfare-identity} is just notation.

\paragraph{Computation.}
With both $F_c^A$ and $F_c^B$ estimated, the expected minimum
cost in an auction with $(n_A, n_B)$ bidders follows directly
from the order-statistic integral. Under the Open regime the
integral runs over both type distributions; under the SME-only
regime it runs only over the SME distribution. Explicitly, under
the Open regime,
\begin{align}
\mathbb{E}[c_{(1)}^{R}] \;=\;
  \int_0^{\overline{c}} \bigl[1 - F_c^A(c)\bigr]^{n_A}
                     \bigl[1 - F_c^B(c)\bigr]^{n_B}\,
   dc \,+\, \underline{c},
\label{eq:expected-min-open}
\end{align}
and under the SME-only regime, with only the $n_A$ SMEs
eligible,
\begin{align}
\mathbb{E}[c_{(1)}^{R'}] \;=\;
  \int_0^{\overline{c}} \bigl[1 - F_c^A(c)\bigr]^{n_A}\, dc
  \,+\, \underline{c}.
\label{eq:expected-min-sme}
\end{align}
The deadweight loss in auction $t$ is the difference between
the two expected minima, scaled by the reference price to put it
in Brazilian reais,
$\text{DWL}_t = \bigl(\mathbb{E}[c_{(1)}^{R'}] -
\mathbb{E}[c_{(1)}^{R}]\bigr) \times P_{\text{ref}}^t$. In
practice, the two integrals are evaluated numerically on a fine
grid via the trapezoidal rule; doing so for each of the 11{,}168
pre-period auctions and aggregating yields the headline welfare
numbers of the next subsection.

\subsection{Aggregate Welfare Results}
\label{sec:welfare-aggregate}

\paragraph{Convite G65 headline.}
Aggregating auction-by-auction across the 11{,}168 Pre-period
Convite G65 auctions that survive filtering, the welfare
decomposition is in Table~\ref{tab:dwl-agg}.

\begin{table}[htbp]
\centering
\caption{Welfare decomposition, Convite G65, 18-month window}
\label{tab:dwl-agg}
\begin{threeparttable}
\small
\setlength{\tabcolsep}{5pt}
\begin{tabular}{lrrr}
\toprule
                                           & R\$ thousand & US\$ thousand & \% of procurement \\
\midrule
Pre-period procurement value (baseline)    & 2{,}500.8    & 714.5          & 100.0\%          \\
\midrule
Buyer's fiscal cost ($\Delta\mathbb{E}[b_{(1)}]$)  & 260.5  & 74.4   & 10.4\%           \\
Winner rent change ($\Delta\mathbb{E}[b-c]$)       & $-34.2$ & $-9.8$ & $-1.4$\%         \\
\midrule
\textbf{Deadweight loss} ($\Delta\mathbb{E}[c_{(1)}]$) & \textbf{294.7} & \textbf{84.2} & \textbf{11.8\%} \\
\bottomrule
\end{tabular}
\begin{tablenotes}
\footnotesize
\item \textit{Notes:} All values are sums over the 11{,}168 Pre-period
Convite G65 auctions with valid ref/bidder composition. USD conversion
at R\$3.50 per USD (the 2016--2019 sample-period average of the BCB
closing rate). Identity
$\Delta\mathbb{E}[c_{(1)}] = \Delta\mathbb{E}[b_{(1)}]
- \Delta\mathbb{E}[b - c]$
holds by construction: 294.7 = 260.5 - $(-34.2)$.
\end{tablenotes}
\end{threeparttable}
\end{table}

The numbers tell a two-sided story. Over the 18-month window, the
SME-only rule on Convite G65 generates a deadweight loss of
\textbf{R\$295k (US\$84k)}, equal to \textbf{11.8\% of procurement
value}. That loss is paid in two currencies. The buyer pays
R\$261k more than it would have otherwise---roughly 88\% of the
total, a straight fiscal cost. The winning SMEs pay another R\$34k
(US\$10k, 12\%) in foregone profit compared with the rents non-SMEs
would have extracted in the open regime. Both sides of the
transaction are worse off, and the loss is not redistributive:
no third party gets to keep what the buyer and the winners both
leave on the table. This is the canonical signature of a
pool-narrowing restriction---and, for a policy often defended on
fiscal-neutrality grounds, a substantive bill.

\paragraph{Decomposition by auction size.}
Does the loss concentrate in small, thin auctions or in larger,
busier ones? Table~\ref{tab:dwl-byN} slices the aggregate by
realized bidder count. The DWL-to-fiscal ratio is remarkably
stable at 111--115\% across every $N$-bin, rising only slightly
with auction size. The intuition is mechanical: a larger auction
typically excludes more non-SMEs in expectation, and each
additional exclusion tightens the remaining book a little more,
so the rent-compression piece grows modestly with $N$. The
headline is that the policy does not just raise prices in a few
thin auctions---the welfare loss is a linear, rather than
outlier-driven, feature of the regime.

\begin{table}[htbp]
\centering
\caption{Welfare decomposition by auction size (Convite G65)}
\label{tab:dwl-byN}
\small
\begin{tabular}{lrrrr}
\toprule
$N$-bin   & Fiscal (R\$k) & DWL (R\$k) & Rent change (R\$k) & DWL/Fiscal \\
\midrule
$N=2$     & 73.5  & 81.9  & $-8.4$   & 111\% \\
$N=3$     & 86.2  & 97.3  & $-11.1$  & 113\% \\
$N=4$     & 62.2  & 70.9  & $-8.7$   & 114\% \\
$N\ge 5$  & 38.6  & 44.6  & $-6.0$   & 115\% \\
\bottomrule
\end{tabular}
\end{table}

\paragraph{Annualization and aggregate scale.}
Annualizing the 18-month Convite figures yields R\$197k/year DWL
(US\$56k/year) for Convite G65 alone. Convite, however, is a small
piece of Group-65 procurement: the modality accounts for roughly
15\% of Group-65 value, while Preg\~ao absorbs 84\% and Dispensa
the last 1\%. If the 11--12\% DWL-to-procurement ratio carries
across modalities---plausible given the cross-modality directional
consistency documented in Section~\ref{sec:robustness}, with the
caveat that per-modality ratios may differ---total Group-65 annual
DWL comes in around \textbf{R\$1.2--1.4M/year (US\$340k--400k)}.
This is for a single product group in a single state.

What about all of Brazilian public procurement? I resist the
temptation to multiply. A credible Brazil-wide aggregate would
require numbers this paper does not observe: the share of the
roughly R\$700 billion annual Brazilian public-procurement
envelope that is both (i) below the R\$80{,}000 SME-only threshold
and (ii) subject to the three-SMEs precondition of LC~123, and
the degree to which the 11.8\% G65-specific ratio transfers to
other sectors. Group-65 medical supplies is unusual in more
than one dimension---high per-item values, oligopolistic
pharmaceutical classes (6531), ANVISA-registered suppliers,
CMED ceiling regulation---and I would not expect a custodial or
office-supply category to deliver the same ratio. A careful
replication across federal (ComprasNet) and municipal procurement,
matched to sector-specific SME-eligibility shares, is on my
research agenda; an unreplicated extrapolation is not in this
paper.

\paragraph{Comparison to \citet{krasnokutskayaseim2011}.}
It is worth asking how a 10--12\% welfare cost compares with what
the structural auction literature has already estimated for softer
versions of the same policy. The closest benchmark is
\citet{krasnokutskayaseim2011}, who study the California 5\%
bid-preference program---a much milder intervention, in which
small firms win if their bid is within 5\% of the lowest
non-favored bid. The California program raises procurement cost
by 0.5\% relative to no preference and shifts the small-firm
share of work from 12.5\% to 15.6\%. Our complete-exclusion regime
produces a 10.4\% fiscal cost and an 11.8\% DWL---roughly
20--25$\times$ their 5\%-preference baseline. That is, of course,
exactly what moving from a 5\% discount to 100\% exclusion should
look like: the intensity of the preference has risen by about
20$\times$, and in asymmetric-bidder models DWL scales
supralinearly in preference intensity \citep[see][Section
IV]{krasnokutskayaseim2011}. The proportionality is useful not
because it validates any particular number but because it tells
us the two estimates live in the same theoretical object. California's
5\% and São Paulo's 100\% are reading off the same welfare curve.

\subsection{The Value-Threshold Exemption Counterfactual}
\label{sec:welfare-vt}

\paragraph{Policy experiment.}
The welfare numbers naturally raise a practical question: can
most of the fiscal cost be recovered without giving up the SME
preference on most items? I consider a \emph{value-threshold
exemption}---keep the SME-only rule on low-value items where the
competition pool is thick and the cost asymmetry is small, but
exempt high-value items where the restriction does the most
damage. A reduced-form approximation on the same data suggests
that exempting the top quartile by reference price recovers
roughly 90\% of the fiscal cost while keeping the preference on
75\% of items \emph{by count}~\citep{genicolo2026a}. The
structural primitives recovered here let me evaluate this design
on its own terms rather than as an extrapolation from a
regression slope.

\paragraph{Procedure.}
For each Pre-period Group-65 Convite auction, I compute both
$E[b|\text{S1}]$ (Open) and $E[b|\text{S3}]$ (SME-only, Post pool).
The fiscal cost of the uniform SME-only rule is
$C_{\text{unif}} = \sum_t (E[b|\text{S3}] - E[b|\text{S1}])
P_{\text{ref}}^t$. For each candidate value threshold $\tau$, I compute
the cost under the threshold rule as
$C(\tau) = \sum_{t: P_{\text{ref}}^t \le \tau}
(E[b|\text{S3}] - E[b|\text{S1}]) P_{\text{ref}}^t$.

\paragraph{Results.}
Table~\ref{tab:cf} reports cost recovery by threshold percentile.

\begin{table}[htbp]
\centering
\caption{Value-threshold exemption counterfactual (Convite G65)}
\label{tab:cf}
\begin{threeparttable}
\small
\begin{tabular}{lrrrr}
\toprule
Threshold percentile & Items exempted & Items kept & Cost recovered &
\% Items kept \\
\midrule
50th (R\$15.43)  &  7{,}204 &  7{,}207 & 98.8\% & 50\% \\
60th (R\$24.77)  &  5{,}764 &  8{,}647 & 97.8\% & 60\% \\
70th (R\$42.88)  &  4{,}322 & 10{,}089 & 96.2\% & 70\% \\
\textbf{75th (R\$57.41)} & \textbf{3{,}600} & \textbf{10{,}811} &
\textbf{95.1\%} & \textbf{75\%} \\
80th (R\$86.00)  &  2{,}882 & 11{,}529 & 93.5\% & 80\% \\
90th (R\$278.00) &  1{,}441 & 12{,}970 & 86.5\% & 90\% \\
\bottomrule
\end{tabular}
\begin{tablenotes}
\footnotesize
\item \textit{Notes:} Value-threshold exemption simulated under the
asymmetric-bidder FPSB model. For auctions with $P_{\text{ref}} > \tau$,
the Open regime (S1) applies; for auctions with $P_{\text{ref}} \le \tau$,
the SME-only regime (S3, Post pool) applies. Cost recovered = fraction
of the uniform-rule fiscal cost eliminated by the exemption.
\end{tablenotes}
\end{threeparttable}
\end{table}

The structural Pareto frontier, plotted in
Figure~\ref{fig:vt-frontier}, is strikingly concave. Exempting
just the top 10\% of items already recovers 87\% of the fiscal
cost; exempting the top 25\% recovers 95\%, leaving the SME-only
rule on three-quarters of items by count. Compared with the 90\%
recovery that the reduced-form approximation delivers at the
same threshold, the structural exercise actually recovers
\emph{more}, because the top-quartile items have (i) higher
reference prices, so their absolute fiscal cost is disproportionate
to their count share, and (ii) higher $N$ on average, so exempting
them does the most to reopen the intensive margin where non-SME
entry would matter most. The practical message for a procurement
regulator is that a small number of high-value items carries most
of the welfare cost of the uniform rule---exempting them is
neither an abandonment of the SME policy nor an ideological
concession, but a direct application of the structural decomposition.
Appendix~H reports the full sensitivity of the Pareto frontier
at twenty threshold percentiles.

\begin{figure}[htbp]
\centering
\includegraphics[width=0.72\textwidth]{fig_v2_vt_frontier.pdf}
\caption{Structural value-threshold exemption Pareto frontier.
   Exempting top-priced items from the SME-only rule recovers
   95\% of the fiscal cost while keeping the rule on 75\% of items.}
\label{fig:vt-frontier}
\end{figure}

%% ============================================================================
\section{Robustness and Extension to Preg\~ao}
\label{sec:robustness}

\paragraph{Preg\~ao extension via \citet{haile2003} bounds.}
Convite is where CPV works, but Convite is only 15\% of
Group-65 procurement by value. The remaining 84\% sits in
Preg\~ao eletr\^onico, an iterative descending reverse auction
with a random-ending-time soft close---structurally closer to
an English auction than to a sealed-bid FPSB, and therefore
outside the domain where CPV point identification applies.
For Preg\~ao I fall back to the \citet{haile2003} bounds
framework, which is weaker but applies to iterative settings.

The relevant inequality is \textbf{H1}: under individual
rationality, no bidder's final bid drops below its own cost, so
the empirical CDF of final bids for each type provides a
pointwise upper bound on that type's cost CDF, $F_c^k(c) \leq
G_b^k(c)$. Applied to the Preg\~ao G65 sample (163{,}004
auctions, 279{,}100 bidder-auctions after ref-price propagation
and filtering), the test flips which type stays invariant
across the break. The non-SME upper bound shifts by $+0.042$ of
reference price (KS $D = 0.061$, vanishing $p$-value in the large
sample), while the SME upper bound is essentially unmoved (mean
shift $+0.0005$, $D = 0.027$). The direction is exactly the
selection story: under SME-only Preg\~ao, the few non-SMEs who
keep bidding despite near-zero win probability are not a random
subsample---they are the non-SMEs with low enough cost or high
enough idiosyncratic reason to bid anyway. Aggregate Preg\~ao
winning-bid shift is $+7.2\%$ of reference price, consistent
with but smaller than the Convite CPV
estimate.\footnote{Restricted apples-to-apples ($N \geq 2$),
mean distinct bidders per G65 auction is $4.12$ in Preg\~ao and
$3.57$ in Convite---Preg\~ao is modestly thicker. The smaller
Preg\~ao effect therefore is not explained primarily by a
thicker pool. The more plausible mechanism is that iterative
open-outcry bidding under the soft-close rule reveals opponent
information during the auction, compressing markups differently
than the sealed-bid Convite mechanism. The Haile-Tamer result
is also a one-sided upper bound rather than a point estimate, so
the underlying cost shift in Preg\~ao may in fact be larger than
the $7.2\%$ the final-bid test can recover.}
Cross-modality validation is therefore directional rather than
point: each modality shows primitive invariance for the type
whose pool is not re-selected by the policy, and each shows a
directionally consistent shift. A type-asymmetric CPV-style
Preg\~ao extension would require a structural model of iterative
bidding under random close, which is on the research agenda but
not in this paper. Appendix~I reports the H1 upper-bound
derivation and the non-SME KS shift disaggregated by $N$-bin.

\paragraph{Alternative SME classifications.}
The SME/non-SME split is the spine of the structural exercise,
so how firms are classified matters, and two reasonable choices
are on the table. The headline specification uses a
firm-registry proxy (\texttt{sme\_proxy} $= 1$ if
\texttt{porte\_empresa} $\in \{01, 03\}$ or \texttt{fornec\_enquad}
$\in \{1, 2\}$; see Appendix~A). A stricter canonical alternative
aggregates the BEC self-declared \texttt{me\_epp} field to the
firm level (winners only, with firm-registry fallback for
losers). Total bidder-auction counts after the CPV clean filter
are nearly identical across the two (18{,}907 proxy vs.\ 18{,}885
canonical), but the Pre-period SME cell---which is the hardest
cell to estimate a density on---shrinks from 2{,}604 to just
372 under canonical. That is a seven-fold reduction in the
data density most critical for identifying $F_A^{\text{Pre}}$
via kernel estimation, and it is the reason the proxy
classification is used for the headline.

What matters for credibility is whether the two classifications
deliver the same economic answer, and they do. Under canonical,
the non-SME distribution remains invariant across the regime
break (KS $D = 0.04$, mean shift $0.013$); the SME distribution
shifts upward (shift $0.009$); and the Pre-period cost gap
between SMEs and non-SMEs widens from essentially zero under
the proxy to $+4.7\%$ under the canonical---in the direction
the policy debate usually assumes (SMEs are costlier). Every
headline sign holds; magnitudes attenuate under the stricter,
noisier classifier. Appendix~C reports the full canonical
comparison, and Appendix~D reports a further tightening that
restricts to \texttt{porte\_empresa == '01'} (strictly micro-firms)
with the same qualitative patterns.

\paragraph{IPV versus common values.}
CPV is built on independent private values (IPV), and if medical
supplies in reality share a common-cost component---input-price
shocks, currency movements, supply-chain disruptions that hit
all bidders at once---the risk-neutral IPV inversion systematically
over-recovers pseudo-costs and biases the decomposition. The
concern is not hypothetical in pharmaceutical procurement, so it
deserves a direct test. Following the logic of
\citet{hendricks2003}, under IPV, after conditioning on bidder
type and auction size, residual dependence of the recovered
$c_{\text{norm}}$ on auction-level signals should vanish. A
regression of $c_{\text{norm}}$ on $\log N$, $\log P_{\text{ref}}$,
bidder type, and a pharma (class 6531) dummy delivers a joint
$F = 1.50$ (2 and 8{,}250 df, clustered at auction level;
$p = 0.22$), with the coefficient on $\log P_{\text{ref}}$ itself
at $+0.0015$ ($p = 0.29$)---statistically and economically
indistinguishable from zero. The IPV restriction is not rejected
at first order. A finer check, the markup-by-$N$ scaling,
deviates from the pure $1/(N-1)$ IPV prediction by 36--55\% at
higher $N$, which is consistent with mild common-values
contamination, with endogenous participation selection, or
with item heterogeneity across $N$-bins. None of these
deviations is large enough to change the structural
decomposition materially. The full diagnostic regression is
reported in Appendix~J.

\paragraph{Boundary-bias sensitivity.}
Kernel density estimators have well-known boundary bias, and
the Silverman bandwidth occasionally generates pseudo-costs
below zero. These artifacts are 7--41\% of the raw sample,
concentrated at $N = 2$ where the bid distribution is thinnest.
Trimming to $\hat c \in (0, 1.5)$ retains 77.7\% of the sample;
a looser trimming at $(-0.5, 2)$ moves the decomposition weights
by less than two percentage points. Appendix~D reports both
variants; nothing in the paper's headline rests on the specific
trim.

\paragraph{Risk aversion.}
The risk-neutral assumption embedded in CPV is not innocent. If
SMEs and non-SMEs face different risk attitudes---a concern in
an environment where small firms may be more exposed to
cash-flow shocks---the pseudo-costs are biased and the
decomposition can absorb that bias asymmetrically. The
\citet{campoguerre2011} extension of GPV to CARA risk aversion
point-identifies the risk parameter from variation in $N$, and
I calibrate a single CARA coefficient $\hat\alpha$ to the
observed markup-by-$N$ scaling reported in
Section~\ref{sec:decomposition}. The fit is tight:
$\hat\alpha \approx 1.25$ matches observed markups to within
half a percentage point across $N \in \{2, 3, 4, 5+\}$, a value
in the range \citet{campoguerre2011} recover for highway
procurement. Under risk-averse recovery, the mean pseudo-cost
drops by $-0.022$ (non-SME) and $-0.026$ (SME) of reference
price---a near-symmetric 4--5\% reduction across types. The
symmetry is what matters: the decomposition weights and the DWL
numbers move by at most two or three percentage points relative
to the risk-neutral baseline, and the primitive-invariance test
results are unchanged. The headline story does not hinge on
assuming risk neutrality. Appendix~K reports the calibration fit
and the by-$N$-bin markup residuals.

%% ============================================================================
\section{Conclusion}
\label{sec:conclusion}

Set-asides for small firms are often evaluated as if they were
equivalent to an accounting tax on public procurement: their fiscal
cost is the gap between what the buyer pays under the restriction and
what the buyer would pay without it. That reading is incomplete. In
S\~ao Paulo state's medical-supply procurement, the 10--11\% price
effect of the SME-only rule runs roughly two-thirds through a channel
the accounting picture does not see: the same small firms bid less
aggressively when their large competitors are excluded. The remaining
third is a shift in who bids, as higher-cost marginal SMEs enter
markets that had been closed to them. The decomposition is not an
abstract exercise: it says that a reform whose only tool is expanding
the SME roster cannot recover most of the efficiency loss, because
the efficiency loss is mechanical, not just compositional. The lever
that does work is narrower: shrinking the \emph{scope} of the rule
itself---exempting the top quartile of items by reference price
recovers 95\% of the fiscal cost while preserving the SME preference
on three-quarters of items.

A step below the policy conclusion, the paper makes two methodological
points that generalize. First, when a policy regime break re-selects
which firms bid but leaves primitive costs untouched, the
type-specific cost distribution whose pool was not re-selected
provides an over-identifying restriction that, while available in
principle to any regime-break setting, is rarely exploited in
applied auction estimation. In Convite, it is the non-SME distribution
that survives the cut cleanly---firms whose outside option of bidding
elsewhere is unaffected by the Group-65 re-designation. Under the
Preg\~ao Haile-Tamer upper bound, it is the SME distribution that
survives. In each modality, the primitive that holds is the one the
policy cannot have moved. I suspect this cross-modality pattern---the
identifying restriction living wherever the pool was not
re-selected---will recur in other settings that exploit regime breaks
in auction data.

Second, the within-auction intensive margin is, under the rigorous
fixed-pool decomposition, the dominant channel of the price effect.
Participation responds dramatically to the policy---non-SME entry
falls 62\% and SME entry rises 52\% across the break---but the
CPV decomposition locates most of the price movement inside the
auction, in how the remaining SMEs bid once their larger
competitors are gone. A calibration that credits every observed
participation shift to the entry channel can be made to invert
the weights, but that calibration double-counts the re-selection
that the CPV intensive margin already captures (Appendix~E). A
rigorous allocation of the price response to endogenous entry
requires the nested fixed-point entry-bidding equilibrium of
\citet{atheyseira2011}, which this paper does not implement and
leaves as natural next work. The headline claim of the paper is
the fixed-pool result; the participation response is a second,
confirmatory fact about the policy's reach.

Three limitations frame the reading. First, the Convite modality
accounts for about 15\% of Group-65 procurement value; the
bulk---Preg\~ao---is addressed by Haile-Tamer upper bounds rather
than point estimates. The directional patterns survive the modality
change, but a full CPV-style analysis of Preg\~ao would require a
structural model of its iterative soft-close format. Second, the
SME classification at the bidder level is a firm-registry proxy;
Appendix~C shows the qualitative conclusions survive the stricter
canonical \texttt{me\_epp} definition, at smaller magnitudes and
smaller effective sample. Third, all welfare numbers are in
18-month, R\$-denominated terms for Convite G65 alone. Extrapolation
to Brazilian public procurement writ large would require inputs---
sector-specific SME-eligibility shares, cross-sector transferability
of the G65 11.8\% DWL ratio---that this paper does not provide.

The \emph{bitter pill} in the Brazilian public-procurement data is
that the restriction delivered what it promised along the
distributional margin---small-firm share of Group-65 procurement
rose sharply---but the unit cost of the work rose with it, and
most of that cost rise happened inside the auction, not at the
point of entry. Each small firm that wins under SME-only does so
against a less competitive book than it faced under the Open
regime, and its bid reflects that. This is the mechanism the
structural model makes visible and the mechanism a procurement
reform has to confront: if the price effect lives in the bid
function, then expanding the SME roster, streamlining registration,
or subsidizing capacity-building will not recover most of the
efficiency loss. The lever that does work is to narrow the
\emph{scope} of the rule---keep the preference on low-value,
SME-rich items, and let competition back in where the restriction
hurts most. That is, I hope, a useful fact for the procurement
officers in Brazil's 5{,}570 municipalities now calibrating Law
14.133/2021, for the roughly 130 jurisdictions worldwide that
operate SME-preference regimes, and for development economists
who work on the recurring trade-off between competition and
distributional policy.

\clearpage
\bibliographystyle{chicago}
\bibliography{References}

\clearpage
\appendix
\section*{Appendix A. SME classification: proxy and canonical variants}

\paragraph{Canonical definition (policy-relevant).}
Under LC~123/2006 (Brazil's Microenterprise Statute), a firm qualifies
for ME (micro-empresa) status if its gross annual revenue does not
exceed R\$360{,}000; for EPP (empresa de pequeno porte), R\$4{,}800{,}000.
The BEC platform captures this classification through a self-declaration
field \texttt{me\_epp} submitted at bid time. This is the
\emph{canonical} SME indicator---the one on which the set-aside rule
legally operates.

\paragraph{Why the canonical field is not directly usable at bidder level.}
In the source BEC microdata tables that feed the paper's bid-level
extraction, \texttt{me\_epp} is recorded only on the winning-bid row
of each purchase-order-item (POI). It is not tied to individual
bidder-auction observations for losing firms. Using the canonical
definition directly is therefore impossible without aggregating across
winning-bid appearances.

\paragraph{Headline proxy.}
The paper's main specification uses a firm-registry proxy:
\[
\text{sme\_proxy} = \mathbf{1}\{\texttt{porte\_empresa} \in \{01, 03\}
\ \text{or}\ \texttt{fornec\_enquad} \in \{1, 2\}\}.
\]
\texttt{porte\_empresa} is the Receita Federal firm-size code (01 = ME
micro, 03 = EPP small, 05 = demais) and \texttt{fornec\_enquad} is the BEC
supplier classification (1 = ME, 2 = EPP). The proxy captures both
ME and EPP classifications as recorded in the firm registry, is constant
across auctions for the same CNPJ, and is available for every bidder.
It yields the largest effective sample for kernel density estimation.

\paragraph{Canonical robustness.}
Appendix~C reports results under a \emph{canonical} classification:
per CNPJ, aggregate \texttt{me\_epp} self-declarations across all
winning-bid appearances in the 2016--2019 sample, and set
$\text{sme\_canonical} = 1$ if the mean self-declaration exceeds 0.5.
For CNPJs that never won (and so have no self-declaration on record),
fall back to \texttt{sme\_proxy}. Agreement with the headline proxy is
approximately 80\%; the most affected cell is the Pre-period SME sample
(2{,}604 firm-auctions under proxy; 372 under canonical). The smaller
canonical SME sample reduces the precision of the Pre SME cost
distribution estimate; every qualitative conclusion survives
(Appendix~C, Tables~C.1--C.2).

\paragraph{Why proxy-as-headline is defensible.}
(i) Statistical: the proxy yields 6--7$\times$ more SME bidder-auctions
in the Pre period, sharpening kernel density estimation at the Pre SME
primitive that our identification strategy relies on most. (ii) Economic:
the proxy's \texttt{porte\_empresa} $\in \{01, 03\}$ captures the same
ME/EPP categorical distinction as LC~123 but through the firm's
\emph{static} registry classification rather than its \emph{bid-time}
self-declaration---this is \emph{more} robust to per-auction
misdeclaration and gaming. (iii) Empirical: the canonical classification
confirms the direction of every headline result (non-SME invariance,
SME shift, decomposition weights, DWL magnitude), and the direction is
robust to the definition.

\section*{Appendix B. Markup by $N$ and the symmetric-GPV benchmark}
\label{app:markup}

Under symmetric IPV-GPV, the equilibrium markup satisfies
$(b - c) \propto 1/(N-1)$, so the ratio of median markups between
consecutive $N$-bins should equal the ratio $(N_1-1)/(N_2-1)$.
Table~\ref{tab:markup-by-n} reports the CPV-recovered markup at each
$N$-bin and compares the observed ratio to the IPV benchmark.

\begin{table}[htbp]
\centering
\caption{Markup by auction size ($N$-bin); IPV benchmark comparison}
\label{tab:markup-by-n}
\small
\begin{tabular}{lrrrr}
\toprule
$N$-bin & Median markup & $n$ (bidder-auctions) &
   $1/(N-1)$ (IPV) & Observed ratio \\
\midrule
$N = 2$     & 0.381 (38\%) & 8{,}431 & 1.00 & 1.00 \\
$N = 3$     & 0.180 (18\%) & 6{,}088 & 0.50 & 0.68 \\
$N = 4$     & 0.120 (12\%) & 3{,}574 & 0.33 & 0.51 \\
$N \ge 5$   & 0.057 (6\%)  & 6{,}253 & 0.20 & 0.29 \\
\bottomrule
\end{tabular}
\end{table}

The observed markup decays in $N$ as the theory predicts, but
about 36--55\% more slowly than the pure IPV prediction at higher
$N$. This moderate over-prediction of competitive pressure is
consistent with (i) mild common-values contamination, (ii)
endogenous participation selection, or (iii) item-class heterogeneity
across $N$-bins (pharmaceuticals concentrated in higher-$N$
auctions). Section~\ref{sec:robustness} formally tests the IPV
null using the \citet{hendricks2003} diagnostic; Appendix~J
reports the full regression.

\section*{Appendix C. Canonical \texttt{me\_epp} robustness}

This appendix re-estimates the structural model with the canonical SME
classification defined in Appendix~A. The canonical flag is the
CNPJ-level aggregation of the BEC \texttt{me\_epp} self-declaration
across all winning-bid appearances, with firm-registry fallback for
never-winners. Table~\ref{tab:canonical-compare} compares the key
quantities.

\begin{table}[htbp]
\centering
\caption{Proxy vs canonical SME classification: headline comparisons}
\label{tab:canonical-compare}
\small
\begin{tabular}{lrrr}
\toprule
Quantity & Proxy (headline) & Canonical & $\Delta$ \\
\midrule
Mean $c/P_{\text{ref}}$, Pre NonSME   & 0.523 & 0.524 & $+0.001$ \\
Mean $c/P_{\text{ref}}$, Pre SME      & 0.520 & 0.548 & $+0.028$ \\
Mean $c/P_{\text{ref}}$, Post NonSME  & 0.524 & 0.537 & $+0.013$ \\
Mean $c/P_{\text{ref}}$, Post SME     & 0.560 & 0.557 & $-0.003$ \\
\midrule
NonSME primitive-invariance shift     & $+0.001$ & $+0.013$ & $+0.012$ \\
SME entry shift                        & $+0.039$ & $+0.009$ & $-0.030$ \\
KS test NonSME ($D$)                  & 0.030 & 0.044 & \\
KS test SME ($D$)                     & 0.081 & 0.075 & \\
\midrule
Cost gap (SME $-$ NonSME, Pre)         & $-0.4\%$ & $+4.7\%$ & \\
Effective SME sample (Pre)             & 2{,}604 & 372 & \\
Effective NonSME sample (Pre)          & 5{,}651 & 7{,}844 & \\
\bottomrule
\end{tabular}
\end{table}

Three observations. First, the non-SME cost distribution remains
economically invariant under both classifications (shifts $<0.02$ of
reference price). The KS $D$ rises from $0.03$ to $0.04$ under canonical,
reflecting the larger noise from the narrower classification, but the
invariance hypothesis is not economically rejected in either. Second,
the entry shift for SMEs attenuates from $+0.039$ to $+0.009$ under
canonical---a factor of $4\times$---because the canonical classification
has 7$\times$ fewer Pre-period SME bidder-auctions and therefore more
noise in the Pre SME distribution. Third, the Pre cost gap between
types sharpens under canonical: SMEs appear $4.7\%$ costlier than
non-SMEs (policy-relevant direction), versus essentially zero under
proxy. The direction---SMEs as higher-cost bidders whose participation
expands under the set-aside---is preserved.

Under canonical, the implied decomposition weights place still
more weight on the intensive margin---$\sim$85\% intensive /
$\sim$15\% entry, versus 66/34 under the headline proxy---because
the smaller Pre SME sample attenuates the entry-margin estimate.
The intensive margin continues to dominate in either case.

The headline specification uses the proxy to retain sample power;
Appendix~C should be read as a directional check on the canonical
classification, not as the preferred point estimate.

\section*{Appendix D. Boundary-bias trimming sensitivity}
\label{app:trim}

Kernel density estimators are subject to well-known boundary bias
when the underlying distribution has compact support. In the CPV
inversion, this manifests as a small fraction of pseudo-costs
falling outside the plausible support $\hat c / P_{\text{ref}} \in
[0, 1]$. Under the headline trim of $(0.001, 1.5)$, the retained
sample is 77.7\% of the pre-filter bidder-auctions and the
weighted decomposition is 66\% intensive / 34\% entry. Under a
more permissive trim of $(-0.5, 2.0)$, the retained sample rises
to 94.3\% and the decomposition shifts by roughly two percentage
points toward the intensive margin. Across trim schemes, the
headline 66/34 split is robust within a narrow two-percentage-point
band. The primitive-invariance test statistic (non-SME $D = 0.03$,
shift $0.001$) is unchanged up to the third decimal under any
trim in $[-0.5, 2.0]$.

\section*{Appendix E. Participation response: reduced-form calibration}
\label{app:endog-calibration}
Section~\ref{sec:endogenous} argues that the 66/34 structural
decomposition is the rigorous result and that a calibration that
credits the observed non-SME exit and SME entry entirely to the
entry channel double-counts re-selection. This appendix reports
the calibration for transparency.

Under the calibration, the CPV intensive margin $\Delta^{\text{int}}$
is retained, but the entry margin $\Delta^{\text{ent}}$ is replaced
by a scaled-up counterpart that absorbs the reduced-form
participation shifts (Table~\ref{tab:entry-reg}): non-SME entry
rates fall $62\%$ and SME entry rates rise $52\%$ conditional on
item class. Allocating these shifts to the entry channel by the
ratio of observed-to-counterfactual participation gives the
by-$N$ weights in Table~\ref{tab:decomp-endogenous}.

\begin{table}[htbp]
\centering
\caption{Decomposition under the participation-response calibration}
\label{tab:decomp-endogenous}
\small
\begin{tabular}{lrrrr}
\toprule
         & \multicolumn{2}{c}{Fixed-pool (headline)} &
           \multicolumn{2}{c}{Participation-calibrated} \\
         & Intensive \% & Entry \% & Intensive \% & Entry \% \\
\midrule
$N = 2$   & 74\% & 26\% & 37\% & 63\% \\
$N = 3$   & 47\% & 53\% & 24\% & 76\% \\
$N = 4$   & 70\% & 30\% & 35\% & 65\% \\
$N \geq 5$& 72\% & 28\% & 36\% & 64\% \\
\midrule
\textit{Average} (unweighted) & $\sim$66\% & $\sim$34\% & $\sim$33\% & $\sim$67\% \\
\bottomrule
\end{tabular}
\end{table}

Two caveats apply. First, the calibration treats every observed
participation shift as entry-channel causation; a fraction of the
shift is equally consistent with the intensive-margin response to
a thinner pool (the remaining bidders face less competition and
rationally bid softer). The CPV fixed-pool decomposition already
incorporates that response; the calibration cannot separate the
two. Second, a proper structural allocation requires the
\citet{atheyseira2011} nested entry-bidding framework with
explicit distributions over entry costs and ex-ante signals,
which this paper does not implement. The 33/67 weights should
therefore be read as an upper bound on the participation share
under a maximally-participation-attributing accounting. The
rigorous decomposition is 66/34, reported in
Section~\ref{sec:decomposition} and Table~\ref{tab:decomp}.

\section*{Appendix F. Derivation of the asymmetric CPV inverse bid function}
\label{app:cpv-derivation}

This appendix derives Equation~\eqref{eq:cpv-inverse}, the
asymmetric-bidder CPV inverse bid function, from the first-order
condition of the type-$k$ bidder's maximization problem. The
derivation follows \citet{campo2003}; I include it here in
self-contained form for the paper to be self-sufficient.

\paragraph{Setup.}
Consider a procurement auction with $n_A$ type-$A$ (SME) bidders
and $n_B$ type-$B$ (non-SME) bidders. Type-$k$ bidders draw
private costs $c \sim F_k$ on $[\underline{c}, \overline{c}]$
with continuous density $f_k$. Bidders submit sealed bids; the
lowest bid wins. Each type-$k$ bidder uses a strictly increasing
equilibrium bid function $\beta_k: [\underline{c}, \overline{c}]
\to \mathbb{R}_+$, yielding the induced equilibrium bid
distribution $G_k(b) = F_k(\beta_k^{-1}(b))$ with density $g_k$.

\paragraph{Type-$k$ bidder's problem.}
A type-$k$ bidder with cost $c$ choosing bid $b$ wins the auction
iff every other bidder submits a higher bid. Under independent
private values, the probability of winning conditional on
opponent counts $(n_A^*, n_B^*)$ (where $n_A^* = n_A - \mathbf{1}\{k=A\}$
and $n_B^* = n_B - \mathbf{1}\{k=B\}$) is
\[
\Pr(\text{win} \mid b) = \bigl[1 - G_A(b)\bigr]^{n_A^*}
\bigl[1 - G_B(b)\bigr]^{n_B^*}.
\]
Expected payoff is $\pi_k(b; c) = (b - c) \cdot \Pr(\text{win}
\mid b)$. Taking the FOC with respect to $b$ and rearranging,
\[
\frac{d \pi_k}{d b} = \Pr(\text{win} \mid b)
- (b - c) \biggl[ n_A^* \frac{g_A(b)}{1 - G_A(b)}
+ n_B^* \frac{g_B(b)}{1 - G_B(b)} \biggr]
\Pr(\text{win} \mid b) = 0.
\]
Since $\Pr(\text{win} \mid b) > 0$ on the equilibrium support,
this simplifies to
\[
1 \;=\; (b - c)
\biggl[ n_A^* \frac{g_A(b)}{1 - G_A(b)}
+ n_B^* \frac{g_B(b)}{1 - G_B(b)} \biggr],
\]
from which
\begin{equation}
c \;=\; b \;-\; \frac{1}{n_A^* \dfrac{g_A(b)}{1 - G_A(b)}
+ n_B^* \dfrac{g_B(b)}{1 - G_B(b)}},
\label{eq:app-cpv-inverse}
\end{equation}
which reproduces Equation~\eqref{eq:cpv-inverse}. The
right-hand side is fully observable: $G_k$ and $g_k$ are
recovered non-parametrically from the empirical distribution of
type-$k$ bids, and the opponent counts $(n_A^*, n_B^*)$ are
observed at the auction level. No parametric assumption on $F_k$
is needed.

\paragraph{Symmetric-within-type special case.}
When $G_A = G_B = G$ and $g_A = g_B = g$, the bracketed hazard
becomes $(n_A^* + n_B^*) g(b)/(1-G(b)) = (N-1) g(b)/(1-G(b))$,
and the formula collapses to $c = b - (1-G(b))/((N-1) g(b))$,
the standard \citet{gpv2000} symmetric procurement GPV.

\paragraph{Identification conditions.}
Two conditions are needed for \eqref{eq:app-cpv-inverse} to
deliver consistent cost primitives:
\begin{itemize}
\item[(i)] \emph{Mixed-type participation}: at least some
auctions must have $(n_A, n_B)$ with both counts positive, so
that both $G_A$ and $G_B$ are non-degenerately identified. In the
Convite G65 sample, 40\% of Pre-period auctions satisfy this.
\item[(ii)] \emph{Monotone bid functions}: $\beta_k$ strictly
increasing in $c$, which the model guarantees under
$F_k$ absolutely continuous. The empirical check in §7
(0\% of pseudo-costs violate $\hat c \leq b$) confirms monotonicity
is not empirically violated.
\end{itemize}

\paragraph{Equilibrium-existence.}
\citet{lebrun1999} and \citet{maskinriley2000} establish
existence and uniqueness of asymmetric FPSB equilibria under
regularity conditions on $F_A, F_B$ that the recovered
distributions satisfy in the Convite sample.

\section*{Appendix G. Bootstrap inference procedure}
\label{app:bootstrap}

Standard errors for the decomposition in
Table~\ref{tab:decomp} are computed via block-bootstrap at the
auction level. This appendix describes the procedure.

\paragraph{Algorithm.}
For $b = 1, \ldots, B = 200$ replicates:
\begin{enumerate}
\item[1.] Draw a simple random sample with replacement of
$n_{\text{auctions}}$ auctions from the Pre-period Convite G65
CPV-eligible sample. Record the replicate's sampled auction IDs
as $\mathcal{B}_b$.
\item[2.] For each auction $t \in \mathcal{B}_b$, include \emph{all}
bidder-auction observations associated with $t$ (this is the
block-bootstrap structure: auctions are the clustering unit,
bidder-auctions are not resampled independently).
\item[3.] Re-estimate the type-specific bid distributions
$\hat G_A^{(b)}, \hat G_B^{(b)}$ and densities $\hat g_A^{(b)},
\hat g_B^{(b)}$ via Gaussian kernel density on the resampled
bids, using Silverman bandwidth computed from the replicate.
\item[4.] Apply the CPV inversion
(Equation~\ref{eq:app-cpv-inverse}) to each bidder-auction in
$\mathcal{B}_b$, recovering pseudo-costs $\hat c_{i,t}^{(b)}$.
\item[5.] Compute the stratum-level decomposition $(\Delta^{\text{int}(b)},
\Delta^{\text{ent}(b)})$ via Equations \eqref{eq:expected-min-open}
and \eqref{eq:expected-min-sme}.
\item[6.] Record $\{\Delta^{\text{int}(b)}, \Delta^{\text{ent}(b)},
\hat\beta_b\}$ for the replicate.
\end{enumerate}

Bootstrap SE of $\Delta^{\text{int}}$ is the sample standard
deviation of $\{\Delta^{\text{int}(b)}\}_{b=1}^{B}$; percentile
95\% CI uses the 2.5th and 97.5th percentiles. Block-bootstrap
at the auction level accommodates the within-auction dependence
of bidder-auction observations (same $n_A, n_B$, same
$P_{\text{ref}}$), which auction-IID bootstrap would not
capture.

\paragraph{Computational complexity.}
$B = 200$ replicates on the 33{,}706-auction Convite G65 sample
takes $\sim$40 minutes on a standard laptop (R/data.table, 8
cores, parallelized across replicates). $B$ was chosen to
balance statistical accuracy against reproducibility: at
$B = 200$, Monte Carlo SE on each replicate's mean is
$< 2\%$ of the point estimate SE.

\paragraph{Bootstrap results.}
Table~\ref{tab:bootstrap-se} reports bootstrap SEs on the
paper's core quantities. The intensive-share bootstrap estimate
(64.2\%; SE 2.5pp; 95\% CI [59.7\%, 69.4\%]) confirms the
headline decomposition's dominance of the within-auction channel;
the ordering ``intensive $>$ entry'' holds in over 99\% of
replicates. Primitive-invariance shifts (non-SME $+0.001$, SME
$+0.039$) retain their economic interpretation: the non-SME
shift's 95\% CI $[-0.008, +0.010]$ excludes any substantive
movement, while the SME shift's CI $[+0.022, +0.051]$ is
cleanly bounded away from zero. Cost-recovery rates at the
50th/75th/90th value-threshold percentiles are tightly
estimated (SE 0.2--1.6pp).

\begin{table}[htbp]
\centering
\caption{Bootstrap standard errors (B = 200 replicates)}
\label{tab:bootstrap-se}
\small
\begin{tabular}{lrrr}
\toprule
Quantity & Point & Boot SE & 95\% CI \\
\midrule
Mean $c/P_{\text{ref}}$, Pre NonSME & 0.523 & 0.004 & $[0.516, 0.530]$ \\
Mean $c/P_{\text{ref}}$, Pre SME    & 0.521 & 0.006 & $[0.510, 0.532]$ \\
Mean $c/P_{\text{ref}}$, Post NonSME & 0.524 & 0.003 & $[0.518, 0.530]$ \\
Mean $c/P_{\text{ref}}$, Post SME   & 0.560 & 0.004 & $[0.551, 0.566]$ \\
\midrule
Non-SME primitive-invariance shift  & $+0.001$ & 0.005 & $[-0.008, +0.010]$ \\
SME entry shift                      & $+0.039$ & 0.007 & $[+0.022, +0.051]$ \\
\midrule
Intensive decomposition share (\%)   & 64.2 & 2.5 & $[59.7, 69.4]$ \\
Entry decomposition share (\%)       & 35.8 & 2.5 & $[30.6, 40.3]$ \\
\midrule
Cost recovery at 50th pctl (\%)     & 98.3 & 0.2 & $[98.0, 98.6]$ \\
Cost recovery at 75th pctl (\%)     & 93.6 & 0.7 & $[92.3, 94.8]$ \\
Cost recovery at 90th pctl (\%)     & 83.5 & 1.6 & $[80.2, 86.4]$ \\
\bottomrule
\end{tabular}
\end{table}

\emph{Note on comparability to main-text headlines:} The
bootstrap is stratum-specific (runs the decomposition on each
resampled auction set without cross-stratum reweighting),
yielding a point estimate of 64.2\% for the intensive share
versus 66.3\% for the $N$-count-weighted figure reported in
Section~\ref{sec:decomposition}. The two differ by less than one
bootstrap standard error and deliver the same economic
conclusion; the main-text weighting is retained for consistency
with the by-$N$-bin presentation.

\section*{Appendix H. Value-threshold counterfactual: full sensitivity}
\label{app:vt-full}

Table~\ref{tab:vt-full} extends the value-threshold counterfactual
reported in Table~\ref{tab:cf} to the full percentile grid from
the 5th to the 95th, allowing a more complete view of the Pareto
frontier between cost recovery and preservation of the SME
preference.

\begin{table}[htbp]
\centering
\caption{Full value-threshold sensitivity (Convite G65)}
\label{tab:vt-full}
\begin{threeparttable}
\small
\begin{tabular}{lrrrr}
\toprule
Threshold \% & Ref price (R\$) & Items exempt & Items keep SME-only &
 Cost recovery \% \\
\midrule
5\%   &  0.26   & 13{,}688 &    723 & 100.0\% \\
10\%  &  0.63   & 12{,}967 &  1{,}444 & 100.0\% \\
20\%  &  2.00   & 11{,}518 &  2{,}893 &  99.9\% \\
30\%  &  4.82   & 10{,}085 &  4{,}326 &  99.8\% \\
40\%  &  9.38   &  8{,}646 &  5{,}765 &  99.4\% \\
50\%  & 15.43   &  7{,}204 &  7{,}207 &  98.8\% \\
60\%  & 24.77   &  5{,}764 &  8{,}647 &  97.8\% \\
70\%  & 42.88   &  4{,}322 & 10{,}089 &  96.2\% \\
\textbf{75\%} & \textbf{57.41} & \textbf{3{,}600} &
 \textbf{10{,}811} & \textbf{95.1\%} \\
80\%  & 86.00   &  2{,}882 & 11{,}529 &  93.5\% \\
85\%  & 144.56  &  2{,}162 & 12{,}249 &  90.9\% \\
90\%  & 278.00  &  1{,}441 & 12{,}970 &  86.5\% \\
95\%  & 767.83  &    721   & 13{,}690 &  77.5\% \\
\bottomrule
\end{tabular}
\begin{tablenotes}
\footnotesize
\item \textit{Notes:} Value-threshold exemption simulated under
the asymmetric-bidder FPSB model. Auctions with $P_{\text{ref}}
> \tau$ revert to the Open regime (S1); auctions with
$P_{\text{ref}} \leq \tau$ stay under SME-only (S3, Post pool).
The Pareto frontier is strongly concave: 95\% of the fiscal cost
is recovered at the 75th percentile threshold ($\tau =$ R\$57.41),
preserving the SME preference on three-quarters of items.
\end{tablenotes}
\end{threeparttable}
\end{table}

Two features of the full frontier are worth noting. First, the
cost-recovery curve is \emph{extremely flat} at low thresholds
(10th--50th percentile), because small items carry negligible
fiscal cost even if many of them are exempted; this is the
concavity that justifies the headline 95\%-recovery-at-25\%-items
trade-off. Second, the curve's slope steepens only past the 85th
percentile, suggesting that the top-15\% value items contain the
bulk of the recoverable fiscal cost.

\section*{Appendix I. Preg\~ao Haile-Tamer bounds: derivation and
by-$N$ detail}
\label{app:ht-pregao}

\paragraph{Haile-Tamer H1 upper bound.}
In an English-auction-like mechanism where participation is
not binding at the close, every bidder has weakly dominant
strategy to remain active as long as the going price exceeds her
cost $c$. Under this individual-rationality condition,
no bidder's final bid drops below her cost: $b^{\text{final}}_{it}
\geq c_{it}$. Consequently, the empirical distribution of final
bids of type $k$ provides a pointwise \emph{upper} bound on the
type-$k$ cost distribution,
\[
F_c^k(c) \;\leq\; G_b^k(c) \qquad \text{for all } c.
\]
This is the \citet{haile2003} H1 bound. It is weaker than point
identification (which would require a structural model of
iterative bidding with random close, beyond the scope of this
paper) but sufficient for the cross-regime KS test that is the
paper's identification check.

\paragraph{By-$N$-bin shifts (non-SME upper bound).}
Table~\ref{tab:ht-byN} reports KS shifts of the non-SME final-bid
distribution across the March 2018 break, disaggregated by
auction size, for the Preg\~ao G65 sample. The aggregate
SME-side KS $D$ of $0.027$ (reported in Section~\ref{sec:robustness},
Preg\~ao paragraph) is economically invariant and not
disaggregated by $N$ here.

\begin{table}[htbp]
\centering
\caption{Preg\~ao Haile-Tamer bounds: non-SME upper-bound shifts by $N$-bin}
\label{tab:ht-byN}
\small
\begin{tabular}{lrrrr}
\toprule
$N$-bin & $n$ Pre & $n$ Post & Shift (\% of ref) & KS $D$ \\
\midrule
$N = 2$   & 7{,}452  & 6{,}704 & $+2.2\%$ & 0.033 \\
$N = 3$   & 6{,}224  & 5{,}664 & $+5.8\%$ & 0.062 \\
$N = 4$   & 4{,}600  & 4{,}165 & $+9.6\%$ & 0.099 \\
$N \geq 5$& 11{,}427 & 9{,}365 & $+10.3\%$ & 0.108 \\
\bottomrule
\end{tabular}
\end{table}

The rising non-SME shift with $N$ is consistent with a selection
story: in larger auctions where the non-SME pool was always
thinner relative to the SME pool, the selection bias on the
remaining post-policy non-SME sub-pool is amplified. The SME
pool, which is not re-selected by the Preg\~ao regime change in
the same way (SMEs continue to be eligible across all auctions),
remains economically invariant across the break.

\section*{Appendix J. IPV vs common-values diagnostic: full regression}
\label{app:ipv-cv}

Following the logic of \citet{hendricks2003}, under IPV the
recovered $c_{\text{norm}}$ should be independent of
auction-level signals conditional on bidder type and auction
size. Table~\ref{tab:ipv-regression} reports the diagnostic
regression.

\begin{table}[htbp]
\centering
\caption{IPV null diagnostic regression}
\label{tab:ipv-regression}
\small
\begin{tabular}{lrr}
\toprule
Regressor & Coefficient & $p$-value \\
\midrule
$\log N$                        & $+0.0094$  & $0.127$ \\
$\log P_{\text{ref}}$           & $+0.0015$  & $0.288$ \\
\midrule
Joint Wald $F(2, 8{,}250)$ on $(\log N, \log P_{\text{ref}})$: &
 \multicolumn{2}{r}{$F = 1.50, p = 0.22$} \\
\bottomrule
\end{tabular}
\end{table}

\emph{Dependent variable}: $c_{\text{norm}}$ (normalized
pseudo-cost). Controls for bidder type and pharma (class 6531)
indicator included but not reported. Standard errors clustered
at the auction level. Sample: 18,907 post-filter bidder-auctions.

The IPV null is not rejected at first order: the joint
hypothesis that $\log N$ and $\log P_{\text{ref}}$ both have zero
coefficient is not rejected ($p = 0.22$). The individual
coefficient on $\log P_{\text{ref}}$, which would reveal a
common reference-price signal in cost, is statistically and
economically indistinguishable from zero. The markup-by-$N$
scaling (Appendix~B) deviates from the $1/(N-1)$ prediction by
36--55\% at higher $N$, but this deviation is compatible with
mild common-values contamination (not rejected by the test) or
with participation selection across $N$-bins.

\section*{Appendix K. Risk aversion: CGPV (2011) calibration}
\label{app:risk-aversion}

\citet{campoguerre2011} extend GPV to allow for constant absolute
risk aversion (CARA) with a common parameter $\alpha$ across
bidders. Under CARA utility $u(x) = -e^{-\alpha x}/\alpha$, the
equilibrium first-order condition becomes
\[
b - c = \frac{1 - e^{-\alpha (b - c)}}{\alpha}
\cdot \frac{1}{N(b) - 1},
\]
where $N(b)$ is the hazard-adjusted opponent count. The
risk-aversion parameter $\alpha$ is point-identified from
variation in $N$.

\paragraph{Calibration procedure.}
I calibrate $\hat\alpha$ to match the observed markup-by-$N$
pattern (Table~\ref{tab:markup-by-n}). Specifically, I search
over $\alpha \in [0, 5]$ for the value that minimizes the
mean-squared-error between observed median markups and the
CARA-implied markups at each of the four $N$-bins.

\begin{table}[htbp]
\centering
\caption{Risk aversion calibration: observed vs CARA-predicted markups}
\label{tab:ra-fit}
\small
\begin{tabular}{lrrr}
\toprule
$N$-bin & Observed markup & CARA-predicted ($\hat\alpha = 1.25$) & Residual \\
\midrule
$N = 2$   & 0.266 & 0.266 & $-0.001$ \\
$N = 3$   & 0.148 & 0.147 & $+0.002$ \\
$N = 4$   & 0.103 & 0.101 & $+0.002$ \\
$N \ge 5$ & 0.058 & 0.062 & $-0.005$ \\
\bottomrule
\end{tabular}
\end{table}

\paragraph{Calibration result.}
$\hat\alpha = 1.25$ matches the observed markups to within 0.5
percentage point across all four $N$-bins. This is in the range
\citet{campoguerre2011} recover for Michigan highway
construction ($\alpha \in [0.8, 1.7]$), suggesting moderate
risk aversion typical of capital-constrained procurement
contractors.

\paragraph{Implied cost adjustment.}
Under risk-averse recovery, the mean pseudo-cost drops by
$-0.0221$ (non-SME) and $-0.0256$ (SME) of reference price---a
near-symmetric $\sim$4\% reduction across types. The key
property is the \emph{symmetry} of the adjustment: both type
distributions shift down by almost the same amount, so the
\emph{relative} position of the two primitives (which is what
the decomposition depends on) is preserved. Consequently:
\begin{itemize}
\item The fixed-pool decomposition moves by at most a few
percentage points and the intensive margin's dominance is
preserved.
\item The primitive-invariance conclusion for non-SMEs is
unaffected---both the risk-neutral and CARA-adjusted estimates
imply economic invariance across the March 2018 break.
\item DWL magnitudes move by less than one percentage point of
procurement value.
\end{itemize}

The headline story does not hinge on assuming risk neutrality;
moderate CARA risk aversion preserves every qualitative conclusion.

\section*{Appendix L. Replication materials}
\label{app:replication}

All scripts, intermediate datasets, and output tables are
archived in the paper's replication package at
\url{https://github.com/darciogm/paper2-me-epp/v2-structural}.
The directory structure is:
\begin{itemize}
\item \texttt{scripts/}: R pipeline (scripts 32--45) and Python
helpers for data extraction.
\item \texttt{data/}: BEC bid-level source parquet files (large;
raw CSV available on request).
\item \texttt{output/tables/}: generated CSVs and \LaTeX{} tables
matching every numerical claim in the paper.
\item \texttt{output/figures/}: generated PDF figures at
publication resolution.
\item \texttt{logs/}: per-script execution logs with timing,
sample sizes, and diagnostic output.
\end{itemize}

Key entry points: script~35 runs the CPV inversion; script~36
builds the decomposition; script~38 produces bootstrap SEs
(Appendix~G); script~39 computes welfare magnitudes. Full
reproduction requires R~4.5+ with \texttt{data.table},
\texttt{arrow}, \texttt{fixest}, \texttt{duckdb}, and
\texttt{ggplot2}. Runtime on a standard laptop: 2--3 hours
end-to-end.

\end{document}
